% Submission version for Computer-Aided Design (Elsevier), built from the same section files as ../p1.tex.
\documentclass[preprint,12pt]{elsarticle}
\usepackage{amsmath,amssymb,amsthm,mathtools}
\usepackage{booktabs,array}
\usepackage{xcolor}
\usepackage[colorlinks=true,linkcolor=blue!50!black,citecolor=blue!50!black,urlcolor=blue!50!black]{hyperref}
\usepackage[margin=1in]{geometry}

% Drafting markers (none should remain in the submission).
\newcommand{\unverified}[1]{\textcolor{red}{[#1 --- not yet checked against the original]}}
\newcommand{\draftnote}[1]{\textcolor{orange!70!black}{[Draft note: #1]}}

\newtheorem{theorem}{Theorem}[section]
\newtheorem{lemma}[theorem]{Lemma}
\newtheorem{proposition}[theorem]{Proposition}
\newtheorem{corollary}[theorem]{Corollary}
\theoremstyle{definition}
\newtheorem{definition}[theorem]{Definition}
\newtheorem{example}[theorem]{Example}
\theoremstyle{remark}
\newtheorem{remark}[theorem]{Remark}

\newcommand{\R}{\mathbb{R}}
\newcommand{\Z}{\mathcal{Z}}
\newcommand{\C}{\mathcal{C}}
\newcommand{\wfs}{\operatorname{wfs}}
\newcommand{\dist}{\operatorname{dist}}
\newcommand{\inter}{\operatorname{int}}
\newcommand{\cl}{\operatorname{cl}}
\newcommand{\conv}{\operatorname{conv}}

\journal{Computer-Aided Design}
\emergencystretch=3em

\begin{document}
\begin{frontmatter}

\title{When does a tolerant B-rep define a solid? Existence, non-existence and uniqueness of realizations}

\author{Guoqing Zhang}
\address{ZENFLUX, Miami Gardens, Florida, USA}

\begin{abstract}
Boundary representations in CAD kernels carry tolerances: edge curves, the boundary curves of adjacent faces and vertex
points do not fit together exactly, and the solid they describe is only implied. We define a realization of a tolerant
B-rep as a PL embedding of its cell complex whose cells stay in their tolerance zones, and study existence,
non-existence and uniqueness. Uniqueness: if every shell has genus at least one, the carriers of the cells are
contractible, and each shell has a product collar containing its zones and a reference surface whose inclusion is a
homotopy equivalence, then the images of all realizations are topologically ambient isotopic, and so are the solids
they bound. The collar condition follows from a distance bound below the weak feature size of a PL reference surface,
for which we give an explicit lower bound in terms of conical radii, with the key lemma machine-checked in Lean.
Existence: a conditional fan-and-cone construction for regular complexes, and a re-intersection construction that
turns an embedded re-trimmed candidate into a realization. Non-existence: a checkable criterion for crossing planar
patches. In OpenCascade 8.0.1, the validity checks decide none of these questions: a sewn wedge passes both checks and
has no realization, and a thick pipe that passes both has two realizations whose solids are not homeomorphic.
\end{abstract}

\begin{keyword}
tolerant boundary representation \sep geometric robustness \sep realization \sep ambient isotopy \sep weak feature
size \sep 3-manifold topology
\end{keyword}

\end{frontmatter}

\section{Introduction}\label{sec:intro}

\subsection{The problem}
A solid in a CAD kernel is stored as a boundary representation: a cell complex (faces, edges, vertices) with
geometry attached to each cell.
\begin{itemize}
\item Each face carries a trimmed patch of a surface.
\item Each edge carries a 3D curve, together with one parameter-space curve (pcurve) for every adjacent face.
\item Each vertex carries a point.
\end{itemize}
In floating point, and after operations such as Booleans, fillets, offsets and sewing, these data do not agree.
The images of the pcurves on the surfaces miss the 3D edge curve, and the edge curves miss the vertex points.
Jackson's local-tolerance approach~\cite{Jackson95}, the published description of Parasolid's, attaches a
tolerance to every vertex, edge and face. The data count as consistent when the deviations lie within the
tolerances. OpenCascade, for instance, requires the pcurve images to lie within the edge tolerance of the 3D curve,
and the edge ends within the vertex tolerance of the vertex point~\cite{OCCT}.

What solid does such a model describe? Qi and Shapiro~\cite{QiShapiro06b} describe the conventional answer: the
representation is valid if there exists an exact solid, with the same boundary combinatorial structure, in the
tolerance zones. This raises two questions:
\begin{enumerate}
\item \emph{Existence.} Is there an embedded surface, with the given cell structure, that respects all the
tolerances?
\item \emph{Uniqueness.} If there are several, do they all bound the same solid up to ambient isotopy?
\end{enumerate}
Downstream algorithms (Booleans, offsets, meshing, analysis) implicitly assume both answers are yes.

\subsection{What we prove}
We call an embedding as in (1) a \emph{realization}: a PL embedding $h$ of the cell complex $K$ with
$h(\inter\sigma)\subseteq\Z(\sigma)$ for every cell $\sigma$, where $\Z(\sigma)$ is the tolerance zone of $\sigma$
(Section~\ref{sec:setting}).

\begin{itemize}
\item \textbf{Uniqueness} (Theorem~\ref{thm:uniqueness}).
  \begin{itemize}
  \item \emph{Hypotheses:} the carriers of the cells are contractible; each shell has a product collar that contains
  its zones and a carried reference surface whose inclusion is a homotopy equivalence; the collars are pairwise
  disjoint; and every shell has genus at least one.
  \item \emph{Conclusion:} the \emph{images} of any two realizations are ambient isotopic, and so are the solids they
  bound.
  \item \emph{Method:} a carrier lemma (a cellular homotopy argument), combined with Waldhausen's theorem that
  incompressible surfaces in a product $F\times I$ can be made horizontal~\cite[Prop.~3.1]{Waldhausen68}.
  \end{itemize}
\item \textbf{The collar condition} (Section~\ref{sec:collar}).
  \begin{itemize}
  \item The collar condition follows from a distance bound below the weak feature size $\wfs(M)$ of a reference
  surface $M$ (Lemma~\ref{lem:collar}). One step uses Moise's approximation of homeomorphisms by PL
  homeomorphisms~\cite[Thm.~35.2]{Moise77}.
  \item For a polyhedral $M$ we prove $\wfs(M)\ge\lambda(M)>0$, where $\lambda$ is defined from the conical radii of
  $M$ (Corollary~\ref{cor:lambda}). The underlying cone lemma is machine-checked in Lean~4 with Mathlib.
  \item We compute lower bounds for $\wfs(M)$ with an octree search, up to stated floating-point qualifications.
  \end{itemize}
\item \textbf{Existence} (Section~\ref{sec:existence}). There are two constructions:
  \begin{itemize}
  \item fans in the normal disks of the edges, with cones at the vertices. This is a conditional theorem for regular
  complexes (Theorem~\ref{thm:lemma0}), proved in Appendix~\ref{sec:proof51}.  Its hypotheses are local conditions on edges and
  vertices, global disjointness conditions (S1), and two non-local topological conditions (V1(b), P0).
  \item re-intersection of the face surfaces, for sharp edges with tight tolerances (Propositions~\ref{prop:B}
  and~\ref{prop:Bprime}; the curved case is proved in Appendix~\ref{sec:proofBprime}). These assume that the re-trimmed
  candidate is embedded and lies in the zones, and turn it into a PL realization.
  \end{itemize}
\item \textbf{Non-existence} (Theorem~\ref{thm:obstruction}). Under explicit conditions, if two planar patches cross
transversally away from the projections (``shadows'') of all edge and vertex zones, there is no realization.
\end{itemize}

\subsection{What the experiments show}
In OpenCascade~8.0.1 (Section~\ref{sec:experiments}):
\begin{itemize}
\item \texttt{BRepCheck\_Analyzer} accepts models however large their tolerances are; in the source, its
tolerance-value check is declared but not implemented.
\item \texttt{BRepAlgoAPI\_Check} gives false alarms on realizable models, and accepts models whose tolerances
exceed the weak feature size of the nominal surface. Among them is a thick pipe with two realizations whose solids
are not homeomorphic: one ties a trefoil into the hole.
\item Sewing a wedge whose faces are slightly displaced produces a model that passes both checks and, by
Theorem~\ref{thm:obstruction}, has no realization. The hypotheses were verified in exact rational arithmetic.
\end{itemize}

\subsection{Related work}

\subsubsection*{Tolerant modelling}
\begin{itemize}
\item Segal~\cite{Segal90} assigned a tolerance to every vertex, edge and face of polyhedral models.
\item Jackson~\cite{Jackson95} gave the local-tolerance B-rep approach.
\item Milenkovic~\cite{Milenkovic88} proved, in the plane and with a uniform $\varepsilon$, that normalized data have an
exact realization within $\varepsilon/2$.
\item Sakkalis, Shen and Patrikalakis~\cite{SSP01} study interval solids: boxes covering the boundary of a known exact
solid, up to homeomorphism. Sakkalis, Peters and Bisceglio~\cite[Thm.~10]{SPB04} prove ambient isotopy for a closed
surface that lies inside a normal tube of a smooth closed surface and meets each normal fibre once. The tube radius
there is bounded by curvature and by a double-normal distance.
\item Qi and Shapiro~\cite{QiShapiro06} develop $\varepsilon$-solidity, which does not require an exact solid to
exist~\cite{QiShapiro06b}.
\item Shapiro~\cite{Shapiro08} gives conditions under which the union of the tolerance zones is homotopy equivalent
to the cell complex. In~\S5.2 (p.~178) he notes that the homotopy condition alone is not sufficient, since a valid
boundary representation must also be an orientable 2-cycle separating space. It is ``not immediately clear'' how to
extend these notions to unions of tolerance zones, though the notion of separation used in isotopy results is likely
to be useful. He also treats the containment of the exact cells in their zones as a much stronger demand.
Realizations satisfy exactly that containment.
\end{itemize}

\subsubsection*{Constructing an exact solid}
\begin{itemize}
\item Andersson, Stewart and Zidani~\cite{ASZ07} and Stewart and Zidani~\cite{SZ08} construct a solid from
inconsistent trimmed-patch data by perturbation (Whitney extension, transfinite interpolation). They \emph{assume}
that the perturbed patches embed:
  \begin{itemize}
  \item ASZ say that adjacent patches need an a posteriori check~\cite[pp.~819--820]{ASZ07};
  \item SZ make a similar assumption~\cite[p.~188]{SZ08}, and derive a normal-vector bound intended to exclude
  extraneous intersections.
  \end{itemize}
\item In an earlier paper~\cite{ASZ06}, the same authors give checkable conditions under which a single composite
parametric surface on a common parameter domain is injective, by a degree argument. They note that separately
trimmed B-rep faces would need extra inside/outside information~\cite[\S3.7]{ASZ06}.
\end{itemize}

\subsubsection*{Isotopy}
\begin{itemize}
\item The weak feature size and the isotopy of offsets below it are studied by Chazal and
Lieutier~\cite{ChazalLieutier07} and by Chazal, Cohen-Steiner and Lieutier~\cite{CCSL09dcg}. The latter
also prove isotopy of offset level sets of general compact sets under Hausdorff
perturbation~\cite[Thms.~4.2,~4.3]{CCSL09cg}.
\item Chazal and Cohen-Steiner~\cite[Thm.~2.1]{CCS05} prove that a surface homeomorphic to $S$, lying in a thickening
of $S$ and separating its sides, is isotopic to $S$. This is close to the topological core of our uniqueness
argument.
\item Andersson, Peters and Stewart~\cite[Thm.~1]{APS00} prove that a one-parameter family of properly joined
curvilinear triangular patches, each injective with non-singular Jacobian and adjacent patches nontangential at every
time, is linked by an ambient isotopy. Their nontangentiality is close to our no-fold clause T in
Proposition~\ref{prop:Bprime}. The settings differ: they assume the perturbed complex is properly joined at every time,
while we start from data whose patches do not meet and construct a realization.
\end{itemize}

\subsection{What is new and what is not}
For each result we say what was known, what we add, and how far we looked. ``We did not find'' refers to the
literature cited in this paper and the searches behind it, not to an exhaustive survey.

\subsubsection*{Uniqueness (Theorem~\ref{thm:uniqueness})}
\begin{itemize}
\item \emph{Known.} The topological core is classical. Waldhausen~\cite[Prop.~3.1]{Waldhausen68} shows that
incompressible surfaces in $F\times I$ can be made horizontal. Chazal and Cohen-Steiner~\cite[Thm.~2.1]{CCS05} prove
that a surface homeomorphic to $S$, lying in a thickening of $S$ and separating its sides, is isotopic to $S$.
Sakkalis, Peters and Bisceglio~\cite[Thm.~10]{SPB04} prove isotopy for a surface in a normal tube of a smooth surface
that meets each fibre once.
\item \emph{What we add.}
  \begin{itemize}
  \item The statement for tolerant B-reps: per-cell zones, and realizations that keep every cell in its own zone.
  \item The carrier lemma (Lemma~\ref{lem:carrier}). With contractible carriers it makes every realization homotopic,
  inside the collar, to the reference surface. Together with the collar hypothesis that the reference surface's
  inclusion is a homotopy equivalence, this gives the $\pi_1$-isomorphism needed for Waldhausen's theorem
  (Lemma~\ref{lem:separation}). So a hypothesis is still needed in place of separation, namely the homotopy-equivalence
  condition, but only for the single reference surface, not for every realization as in~\cite{CCS05}, and no realization needs to meet each fibre once as
  in~\cite{SPB04}.
  \item A collar from the weak feature size (Lemma~\ref{lem:collar}), for PL reference surfaces. These may have sharp
  edges and triangulated small fillets. The result of~\cite{SPB04} is stated for a smooth closed reference surface, with
  the tube radius bounded by curvature and by a double-normal distance. Each result has its own setting restriction:
  ours needs a PL reference.
  \end{itemize}
\item \emph{Limits.}
  \begin{itemize}
  \item The collar condition is a hypothesis. Lemma~\ref{lem:collar} supplies it only when the reference surface is a
  carried PL embedding of the shell, every zone lies in $N_d(M)$ with $d<r<\wfs(M)$, the closed $r$-neighbourhoods of
  different shells are pairwise disjoint, and the genus is at least one.
  \item Shells of genus $0$ are excluded.
  \item The isotopy is topological, and of images, not of embeddings.
  \end{itemize}
\end{itemize}

\subsubsection*{The explicit bound $\wfs(M)\ge\lambda(M)$ (Corollary~\ref{cor:lambda})}
\begin{itemize}
\item \emph{Known.} That $\wfs(M)>0$ for compact semialgebraic $M$. Fu~\cite[\S5.3]{Fu85}
proves that the distance function of a compact semialgebraic set has finitely many critical values. di Rocco, Eklund and
Gäfvert~\cite[Thm.~2.5]{dREG22} state positivity and attribute it to Fu. Chazal and
Lieutier~\cite[Prop.~3.6]{ChazalLieutier07} state finiteness for bounded subanalytic open sets, attribute it to Fu for
semialgebraic sets, and refer to an earlier paper of theirs for the adaptation. Fu defines semialgebraic sets by one
finite system of inequalities; his argument (Tarski--Seidenberg, a finite decomposition into smooth submanifolds, his
Lemma~4.3 and Sard's theorem) does not use that restriction; that it covers polyhedral surfaces is our reading.
\item \emph{What we add.} An explicit lower bound in terms of conical radii, whose key cone lemma is machine-checked in
Lean~4; the positivity $\lambda(M)>0$ is proved on paper and not formalized. We did not find such a bound. That does not
make it new.
\end{itemize}

\subsubsection*{Existence (Theorem~\ref{thm:lemma0}, Propositions~\ref{prop:B} and~\ref{prop:Bprime})}
\begin{itemize}
\item \emph{Known.}
  \begin{itemize}
  \item Milenkovic~\cite{Milenkovic88}: existence in the plane, with one uniform tolerance.
  \item Andersson, Stewart and Zidani~\cite{ASZ06}: checkable injectivity conditions for one composite patch on a common
  parameter domain.
  \item \cite{ASZ07,SZ08}: they construct a solid from inconsistent data but assume that the perturbed patches embed.
  \item Andersson, Peters and Stewart~\cite[Thm.~1]{APS00}: ambient isotopy along a family that stays properly joined.
  Their nontangentiality is close to our no-fold clause T.
  \end{itemize}
\item \emph{What we add.} Sufficient conditions under which a 3D B-rep with per-cell tolerance zones has a realization
of the whole cell complex. They are of two different kinds.
  \begin{itemize}
  \item Theorem~\ref{thm:lemma0} constructs the realization from conditions on the data. It is proved for regular
  complexes only, and it is conditional: two of its hypotheses are non-local.
  \item Propositions~\ref{prop:B} and~\ref{prop:Bprime}, for sharp edges with tight tolerances, assume that a
  re-trimmed candidate surface is already embedded and lies in the zones (in Proposition~\ref{prop:Bprime}, clauses C
  and B1--B3), and show that it gives a realization; for curved faces this needs a PL approximation. They check a given
  candidate rather than derive embeddedness. Proposition~\ref{prop:B} is for planar faces and vertices of degree three;
  Proposition~\ref{prop:Bprime} needs $C^2$ regular face parametrizations, without poles or cone apices. On curved models these clauses are
  checked on samples in double precision, and the vertex part of clause T is not checked.
  \end{itemize}
\item \emph{How far we looked.} We did not find existence results of either kind for 3D B-reps with per-cell
tolerance zones.
\item \emph{Not covered.} Sewn models with tangent edges and OpenCascade's vertex tolerances. Whether they have
realizations is open (Section~\ref{sec:open}).
\end{itemize}

\subsubsection*{Non-existence (Theorem~\ref{thm:obstruction})}
\begin{itemize}
\item \emph{What we add.} A checkable criterion for planar patches. We apply it to sewn models that pass both of
OpenCascade's checks.
\item \emph{Limits.} The theorem concerns one pair of planar faces whose closed cells are disks bounded by single wires,
with simply connected edge zones (O3). A version for curved faces is open.
\item \emph{How far we looked.} We did not find a non-existence criterion for tolerant B-reps in the literature we
read.
\end{itemize}

\subsubsection*{Experiments (Section~\ref{sec:experiments})}
\begin{itemize}
\item \emph{What we add.} The observation that OpenCascade's validity checks decide neither existence nor uniqueness.
\item \emph{Scope.} The evidence is a handful of parts in one kernel version (8.0.1). We did not test ACIS or
Parasolid.
\end{itemize}

\section{Tolerant B-reps and realizations}\label{sec:setting}

\subsection{Complexes}
$K$ is a finite cell complex in the B-rep sense. Its 0-cells are points, its edges are arcs or closed curves, and
its faces are compact surfaces, possibly with holes. The open face is the face minus all its edges, seams included. The underlying space $|K|$ is a
disjoint union of closed connected surfaces $|K_1|,\dots,|K_n|$ (the \emph{shells}). We write $\tau\le\sigma$ when $\tau$ is a face of $\sigma$ (including
$\tau=\sigma$). For the uniqueness theorem we assume $K$ is \emph{regular}: every closed cell $\bar\sigma$ is a
ball whose boundary is a subcomplex. B-rep complexes with seam edges, closed edges or faces with holes are not
regular, and must first be subdivided. A new cell inherits the zone of the open cell of $K$ that contains it.
This keeps the zones monotone, and a realization of the original model is a realization of the subdivision, since
$h(\inter\sigma')\subseteq h(\inter\sigma)\subseteq\Z(\sigma)$ for a new cell $\sigma'\subseteq\sigma$. The carrier
condition of Section~\ref{sec:uniqueness} must then hold for the subdivided complex. A new edge inside a face $f$
has carrier $\Z(f)$, so subdividing faces adds the hypothesis $\pi_1(\Z(f))=0$; contractible zones satisfy it. A map of $|K|$ is \emph{PL} if it is linear on the simplices of some triangulation of $|K|$ that subdivides its
cells. In dimension two, any two such triangulations are related by a cell-preserving PL homeomorphism (standard; not
separately checked against a source), so the notion does not depend on the choice. New vertices created inside an edge or face inherit its zone, so the carrier condition requires those
zones to be path-connected. The existence results say where they apply to B-rep
complexes directly.

\subsection{Geometric data and zones}
\begin{itemize}
\item Each vertex $v$ has a point $p_v$ and a tolerance $\tau_v$.
\item Each edge $e$ has a $C^2$ curve $\gamma_e\colon[0,1]\to\R^3$ and a tolerance $\tau_e$.
\item Each face $f$ has a $C^2$ surface $S_f$, a compact trimmed patch $P_f\subseteq S_f$, and a tolerance
$\tau_f$. The boundary of $P_f$ consists of the images $\beta_{f,e}=S_f\circ(\text{pcurve})$ of the pcurves of the
edges $e<f$.
\end{itemize}
The \emph{deviation} of $e$ is $\mathrm{dev}_e=\max_{f>e}\max_t|\beta_{f,e}(t)-\gamma_e(t)|$, measured at equal
parameters. OpenCascade's SameParameter property asks for $\mathrm{dev}_e\le\tau_e$ in this sense; its format
specification defines the edge tolerance as a one-sided distance from the pcurve images to the 3D
curve~\cite{OCCT}.

\begin{definition}[Zones]\label{def:zones}
We use \emph{augmented} open zones:
\begin{itemize}
\item $\Z(v)=B(p_v,\tau_v)$;
\item $\Z(e)=N_{\tau_e}(\gamma_e)\cup\bigcup_{v<e}\Z(v)$;
\item $\Z(f)=N_{\tau_f}(P_f)\cup\bigcup_{e<f}\Z(e)$.
\end{itemize}
Here $N_r$ is the open $r$-neighbourhood. Augmentation makes $\Z$ monotone: $\tau\le\sigma$ implies
$\Z(\tau)\subseteq\Z(\sigma)$.
\end{definition}

\begin{definition}[Realization]\label{def:realization}
A \emph{realization} of the tolerant B-rep $T=(K,\text{data},\tau)$ is a PL embedding $h\colon|K|\to\R^3$ with
$h(\inter\sigma)\subseteq\Z(\sigma)$ for every cell $\sigma$.
\end{definition}

\begin{remark}\label{rem:opencell}
\begin{itemize}
\item \emph{Open-cell condition.} We require only $h(\inter\sigma)\subseteq\Z(\sigma)$. The closed-cell
condition with unaugmented neighbourhoods, $h(\bar\sigma)\subseteq N_{\tau_\sigma}(\cdot)$, usually fails for every
candidate, mainly because face tolerances are tiny: $h(v)$ would have to lie within $\tau_f$ of every incident patch. With augmented zones the open-cell
condition implies $h(\bar\sigma)\subseteq\Z(\sigma)$, since $\bar\sigma$ is the union of the open cells
$\tau\le\sigma$.
\item \emph{Tameness.} Requiring $h$ to be PL rules out wild embeddings such as Alexander's horned sphere.
\end{itemize}
\end{remark}

\begin{definition}[Carriers]
The \emph{carrier} of a cell is $\C(\sigma)=\bigcup_{\tau\le\sigma}\Z(\tau)$. For augmented zones
$\C(\sigma)=\Z(\sigma)$. A continuous map $g\colon|K|\to\R^3$ is \emph{carried} if $g(\bar\sigma)\subseteq
\C(\sigma)$ for every $\sigma$. By Remark~\ref{rem:opencell}, every realization is carried.
\end{definition}

\begin{definition}[Solid]\label{def:solid}
Let $S_1,\dots,S_n\subset\R^3$ be pairwise disjoint closed surfaces. The \emph{depth} of a point $x\notin\bigcup S_j$
is the number of $j$ for which $x$ lies in the bounded component of $\R^3\setminus S_j$. The solid bounded by
the $S_j$ is the closure of the set of points of odd depth.
\end{definition}
This matches the usual B-rep convention of outer shells and cavities when the shells are consistently oriented;
a consistency check of the orientation is separate.

\section{Uniqueness for shells of positive genus}\label{sec:uniqueness}

\subsection{Hypotheses}
\begin{itemize}
\item \textbf{Carrier condition.} $\C(v)$ is path-connected for every vertex $v$, and $\C(\sigma)$ is path-connected
with $\pi_k(\C(\sigma))=0$ at every base point, for every $k$-cell $\sigma$ with $k\ge1$. Contractible carriers suffice.
\item \textbf{Collar condition.} For each shell $K_i$ there are a compact PL submanifold $U_i\subset\R^3$, a
closed surface $M_i$, a PL homeomorphism $\Psi_i\colon U_i\to M_i\times[-1,1]$, and a carried PL embedding
$\varphi_i\colon|K_i|\to\R^3$ (a \emph{reference}), with these properties:
  \begin{itemize}
  \item all zones of the cells of $K_i$ lie in $\inter U_i$;
  \item $\varphi_i(|K_i|)\subseteq\inter U_i$, and the inclusion $\varphi_i(|K_i|)\hookrightarrow U_i$ is a homotopy
  equivalence;
  \item the $U_i$ are pairwise disjoint.
  \end{itemize}
  (The reference need not be a level $\Psi_i^{-1}(M_i\times\{t\})$. This weaker form is what
  Lemma~\ref{lem:collar} can deliver.) Note that $M_i$ is connected, since $M_i\simeq U_i\simeq|K_i|$; it is
  orientable, since $M_i\times I\cong U_i\subset\R^3$; and it is not a sphere when $|K_i|$ is not, since then $\pi_1(M_i)\cong\pi_1(|K_i|)\ne1$.
\item \textbf{Genus.} Each $|K_i|$ has genus at least one.
\end{itemize}

\subsection{Lemmas}

\begin{lemma}[Carrier lemma]\label{lem:carrier}
Under the carrier condition, any two carried maps $f,g\colon|K|\to\R^3$ are homotopic by a homotopy $H$ with
$H(\bar\sigma\times I)\subseteq\C(\sigma)$ for every $\sigma$.
\end{lemma}
\begin{proof}
Define $H$ on $|K|\times\{0,1\}\cup|K^{(k)}|\times I$ by induction on $k$, keeping $H(\bar\sigma\times I)
\subseteq\C(\sigma)$, with $H(\cdot,0)=f$ and $H(\cdot,1)=g$.
\begin{itemize}
\item $k=0$: for a vertex $v$, the points $f(v),g(v)$ lie in the path-connected set $\C(v)$; join them by a
path in $\C(v)$.
\item $k\ge1$: let $\sigma$ be a $k$-cell. $\bar\sigma\times I$ is a $(k+1)$-ball whose boundary
$\bar\sigma\times\{0,1\}\cup\partial\bar\sigma\times I$ is a $k$-sphere, and $H$ is already defined on it:
  \begin{itemize}
  \item on $\bar\sigma\times\{0,1\}$ it is $f$ or $g$, with image in $\C(\sigma)$;
  \item on $\partial\bar\sigma\times I$ it is defined by induction, with image in
  $\bigcup_{\tau<\sigma}\C(\tau)\subseteq\C(\sigma)$.
  \end{itemize}
  Since $\pi_k(\C(\sigma))=0$, $H$ extends over $\bar\sigma\times I$ with image in $\C(\sigma)$.
\end{itemize}
The pieces agree on common boundaries, and a finite union of closed sets glued continuously gives a continuous
$H$.
\end{proof}

\begin{lemma}[Separation]\label{lem:separation}
Assume the carrier and collar conditions, and let $h$ be a realization. Then $G=h(|K_i|)$ is a closed surface in
$\inter U_i$, and the inclusion $G\hookrightarrow U_i$ is a homotopy equivalence.
\end{lemma}
\begin{proof}
\begin{itemize}
\item By Lemma~\ref{lem:carrier} on $K_i$, $h$ and $\varphi_i$ are homotopic through maps with image in
$\bigcup\C(\sigma)\subseteq\inter U_i$, so the homotopy takes place in $U_i$.
\item $\varphi_i\colon|K_i|\to U_i$ is a homeomorphism onto $\varphi_i(|K_i|)$ followed by the inclusion, which is a
homotopy equivalence by the collar condition. Hence so is $h\colon|K_i|\to U_i$.
\item Since $h$ is a homeomorphism onto $G$, the inclusion $G\hookrightarrow U_i$, which equals $h\circ h^{-1}$,
is a homotopy equivalence.
\end{itemize}
\end{proof}

\begin{lemma}[Horizontality]\label{lem:horizontal}
Let $U\cong_{\Psi}M\times[-1,1]$ be as in the collar condition, with $M$ not a sphere, and let $G\subset\inter U$
be a closed PL surface whose inclusion into $U$ is a homotopy equivalence. Then there is an ambient isotopy of
$\R^3$, supported in $U$ and fixed on $\partial U$, that carries $G$ to $\Psi^{-1}(M\times\{0\})$.
\end{lemma}
\begin{proof}
\begin{enumerate}
\item[(a)] \emph{$G$ is incompressible.}
  \begin{itemize}
  \item By hypothesis, $\pi_1(G)\to\pi_1(U)$ is an isomorphism.
  \item The boundary of a compressing disk would be a curve that is essential in $G$ and null-homotopic in $U$,
  which this rules out. That is incompressibility in the sense of Waldhausen~\cite[p.~58]{Waldhausen68}.
  \item $G\simeq U\simeq M$, so $G$ is not a sphere and case (b) of that definition does not arise. We do not use the
  converse direction, which needs the loop theorem.
  \item $G$ is a closed orientable surface in an orientable 3-manifold, hence two-sided, and it is disjoint from
  $\partial U$; so it satisfies Waldhausen's conventions~\cite[\S0]{Waldhausen68}.
  \end{itemize}
\item[(b)] \emph{$G$ can be made horizontal.}
  \begin{itemize}
  \item Apply Waldhausen's Proposition~3.1~\cite[pp.~65--68]{Waldhausen68} to $U\cong M\times I$ ($M$ orientable,
  not a sphere) and the incompressible surface $G$. Since $\partial G=\emptyset$, the condition
  ``$\partial G\subseteq F\times1$'' holds vacuously.
  \item It gives an isotopy, constant on $\partial U$, from $G$ to a surface $G'$ whose projection
  $p\colon G'\to M$ is a homeomorphism onto its image.
  \item $p(G')$ is open (invariance of domain) and closed in the connected surface $M$, so $p(G')=M$. Hence $G'$ is
  the graph of a function $s\colon M\to(-1,1)$.
  \item Waldhausen's proof of Proposition~3.1 uses that $F\times I$ is irreducible, explicitly in the base case
  ($F$ a disc, annulus or pair of pants, p.~66) and, through the induction on p.~67, for compact surfaces $F'$ with
  boundary obtained by cutting. Lemma~\ref{lem:irreducible} below covers all these cases.
  \end{itemize}
\item[(c)] \emph{Flatten.} The PL homeomorphisms of the fibres $\{x\}\times[-1,1]$ that fix both ends and send $s(x)$
to $(1-t)s(x)$, $t\in[0,1]$, give an ambient isotopy of $U$, fixed on $\partial U$, from $G'$ to
$\Psi^{-1}(M\times\{0\})$. This isotopy is topological.
\item[(d)] The isotopy of (b) is already ambient: in Waldhausen's paper an isotopy of subspaces $Z_1,Z_2$ of $X$ is by
definition an isotopy of the identity map of $X$ carrying $Z_1$ to $Z_2$~\cite[p.~57]{Waldhausen68}, and Proposition~3.1
gives one constant on $\partial U$. Both isotopies are fixed on $\partial U$ and extend by the identity to $\R^3$.
\end{enumerate}
\end{proof}

\begin{lemma}\label{lem:irreducible}
Let $F$ be a compact orientable PL surface other than the 2-sphere. Then $F\times I$ is irreducible: every PL
2-sphere $S\subset\inter(F\times I)$ bounds a PL 3-cell in $F\times I$. (Following Waldhausen's convention that surfaces
are properly embedded, a 2-sphere in $F\times I$ is disjoint from the boundary.)
\end{lemma}
\begin{proof}
\emph{Case $\partial F\ne\emptyset$.}
\begin{itemize}
\item Then $F\times I$ is a handlebody: a regular neighbourhood $H$ of a finite graph in $\R^3$, with $\partial H$
connected (standard; not separately checked against a source). We work with $H\subset\R^3$.
\item By the PL Schoenflies theorem~\cite[Thm.~17.12]{Moise77}, $S$ bounds a PL 3-cell $B\subset\R^3$. The set
$\R^3\setminus H$ is connected (Alexander duality, as $\partial H$ is connected), unbounded and disjoint from $S$, so it
lies in the unbounded component of $\R^3\setminus S$. Hence $B\subseteq H$.
\end{itemize}
\emph{Case $F$ closed, of genus at least one.}
\begin{itemize}
\item Let $q\colon\tilde F\times I\to F\times I$ be the universal covering, with deck group $\Gamma$ acting by
$\gamma\times\mathrm{id}$. $\tilde F$ is a simply connected non-compact PL surface without boundary, so it is PL
homeomorphic to $\R^2$ (standard; not separately checked against a source). Hence $\tilde F\times I\cong_{PL}\R^2\times
I\subset\R^3$.
\item $S$ is simply connected, so it lifts to a PL 2-sphere $\tilde S\subset\R^2\times(0,1)$, and the lifts $\gamma\tilde S$,
$\gamma\in\Gamma$, are pairwise disjoint. By the PL Schoenflies theorem~\cite[Thm.~17.12]{Moise77}, $\tilde S$ bounds a
PL 3-cell $\tilde B\subset\R^3$. Points of $\R^3$ outside $\R^2\times[0,1]$ are joined to infinity in the complement of
$\tilde S$, so $\tilde B\subset\R^2\times I$.
\item Let $\gamma\ne1$. The spheres $\tilde S$ and $\gamma\tilde S$ are disjoint, so the cells $\tilde B$ and
$\gamma\tilde B$ are either disjoint or nested. If $\gamma\tilde B\subseteq\tilde B$, then $\gamma^n\tilde B\subseteq\tilde
B$ for all $n\ge1$. $\Gamma$ is torsion-free, so the $\gamma^n$ are distinct, and infinitely many deck transformations
move the compact set $\tilde B$ to meet itself. That contradicts proper discontinuity. The case $\tilde B\subseteq
\gamma\tilde B$ is the same with $\gamma^{-1}$.
\item So the translates of $\tilde B$ are pairwise disjoint, $q$ is injective on $\tilde B$, and $q(\tilde B)$ is a PL
3-cell in $F\times I$ with boundary $S$.
\end{itemize}
Also standard and not separately checked: that $\pi_1$ of a closed surface is torsion-free, and that deck groups act
properly discontinuously.
\end{proof}

\subsection{The theorem}

\begin{theorem}[Uniqueness]\label{thm:uniqueness}
Assume the carrier condition, the collar condition, and that every shell has genus at least one. Then:
\begin{enumerate}
\item the image of any realization is ambient isotopic (by a topological isotopy) to
$\bigcup_i\Psi_i^{-1}(M_i\times\{0\})$;
\item the images of any two realizations are ambient isotopic;
\item the solids they bound (Definition~\ref{def:solid}) are ambient isotopic.
\end{enumerate}
\end{theorem}
\begin{proof}
\begin{itemize}
\item By Lemma~\ref{lem:separation}, each $G_i=h(|K_i|)$ satisfies the hypotheses of Lemma~\ref{lem:horizontal} in
$U_i$. Apply that lemma shell by shell. The supports $U_i$ are pairwise disjoint, so the isotopies $\Phi^i$ can be
run simultaneously; their end map is a homeomorphism $\Phi$ of $\R^3$.
\item $\Phi(G_j)=\Psi_j^{-1}(M_j\times\{0\})$ for every $j$: on $U_j$, every $\Phi^i$ with $i\ne j$ is the identity,
because its support $U_i$ is disjoint from $U_j$.
\item By the Jordan--Brouwer separation theorem each closed surface separates $\R^3$ into a bounded and an unbounded
component. A homeomorphism of $\R^3$ maps sets with compact closure to sets with compact closure, so $\Phi$ maps the
bounded component of $\R^3\setminus G_j$ onto the bounded component of $\R^3\setminus\Phi(G_j)$. So depth is
preserved, and $\Phi$ maps the solid bounded by the $G_j$ onto the solid bounded by the levels.
\item For two realizations with maps $\Phi$ and $\Phi'$, use $\Phi'^{-1}\circ\Phi$.
\end{itemize}
\end{proof}

\begin{remark}
\begin{itemize}
\item \emph{What is not covered.} Spheres (genus 0) are excluded, and the isotopy is of images, not of
embeddings.
\item \emph{Relation to Chazal--Cohen-Steiner.} Chazal and Cohen-Steiner~\cite[Thm.~2.1]{CCS05} prove that a compact
orientable surface $S'\cong S$ lying in a thickening of $S$ and separating its sides is isotopic to $S$ in the
thickening. Lemmas~\ref{lem:separation} and~\ref{lem:horizontal} play the same role here. What is specific to
tolerant B-reps is the set-up: per-cell zones, cell-respecting realizations, the carrier lemma, and $\pi_1$ in place
of separation, which the carrier lemma provides.
\item \emph{The collar condition cannot be dropped.} The thick pipe of Section~\ref{sec:experiments} passes both
OpenCascade checks and has two realizations whose solids are not even homeomorphic.
\item \emph{Standard facts used without separate checking against a source:} continuity of cell-by-cell extension,
orientability of closed surfaces in $\R^3$, invariance of domain, the Jordan--Brouwer separation theorem, and that
the fibrewise flattening is an ambient isotopy.
\end{itemize}
\end{remark}

\section{From the weak feature size to the collar condition}\label{sec:collar}

\subsection{Weak feature size}
Let $M\subset\R^3$ be compact, $d(p)=\dist(p,M)$, and let $\Gamma(p)$ be the set of nearest points of $p$ in $M$.
A point $p\notin M$ is \emph{critical} if $0\in\conv\{p-x:x\in\Gamma(p)\}$. The \emph{weak feature size}
$\wfs(M)$ is the infimum of the positive critical values of $d$~\cite[\S2.2]{CCSL09dcg}; see also
\cite{ChazalLieutier07}. Unlike the reach, it
stays positive at sharp edges and small fillets. A surface with a concave edge has reach $0$, and a fillet of
radius $r$ caps the reach at $r$.

\begin{lemma}[Collar]\label{lem:collar}
Let $M$ be a closed connected PL surface in $\R^3$, and let $0<d<r<\wfs(M)$.
\begin{enumerate}
\item[(a)] The closed neighbourhood $\bar N_r(M)$ is homeomorphic to $M\times[-1,1]$, with $M$ going to
$M\times\{0\}$.
\item[(a$'$)] There is a compact PL submanifold $U$ with $\bar N_d(M)\subseteq\inter U$, $U\subseteq\inter\bar
N_r(M)$ and $U\cong_{PL}M\times[-1,1]$.
\item[(b)] Suppose $M$ has genus at least one, and let $U$ be as in (a$'$). Suppose further that $M=\varphi(|K_i|)$ for a carried PL embedding $\varphi$ of a
shell, that every zone of a cell of that shell lies in $N_d(M)$, and that the sets $\bar N_r(M_i)$ of different shells
are pairwise disjoint. Then the collar condition of Section~\ref{sec:uniqueness} holds, with $U_i=U$.
\end{enumerate}
\end{lemma}
\begin{proof}
Write $d=d_M$. By the Jordan--Brouwer separation theorem, $\R^3\setminus M$ has two components, a bounded one
$\Omega_-$ and an unbounded one $\Omega_+$.

(a) \emph{Step 1: small radii.}
\begin{itemize}
\item Work in a large simplicial cube $Q$ containing $\bar N_r(M)$, triangulated so that $M$ is a subcomplex. The
restriction of $d$ to $Q$ is semialgebraic (Tarski--Seidenberg), proper, nonnegative, with zero set the compact set
$M$. Durfee's definition of a (semialgebraic) \emph{rug function}~\cite[Def.~3.1]{Durfee83} asks exactly this and no
smoothness, so $d|_Q$ is a rug function for $M$ in $Q$. For suitably small $\delta>0$, the set $\bar N_\delta(M)=d^{-1}[0,\delta]$
is therefore a semialgebraic neighbourhood of $M$ in $Q$~\cite[Def.~3.4]{Durfee83}.
\item A PL regular neighbourhood $R$ of $M$ in $Q$ is also a semialgebraic neighbourhood~\cite[Prop.~3.7]{Durfee83}.
By the uniqueness of semialgebraic neighbourhoods~\cite[Prop.~3.5]{Durfee83} there are homeomorphisms $h_t\colon Q\to Q$
with $h_t|_M=\mathrm{id}$ for all $t$ and $h_1(\bar N_\delta(M))=R$. So $h_1$ is a homeomorphism of pairs
$(\bar N_\delta(M),M)\to(R,M)$ that fixes $M$ pointwise. We take $\delta<r$.
\item $M$ is a closed orientable surface in $\R^3$, hence two-sided, so $R\cong M\times[-1,1]$ with $M$ going to
$M\times\{0\}$. Hence $(\bar N_\delta(M),M)\cong(M\times[-1,1],M\times\{0\})$. Its frontier is $S_\delta=\{d=\delta\}$,
because $\delta$ is not a critical value, so every point with $d=\delta$ is a limit of points with $d>\delta$ (the
directional-derivative argument of Step~3). $S_\delta$ has two
components, the images of $M\times\{\pm1\}$, and they lie in different components of $\R^3\setminus M$. The two halves
$M\times(0,1]$ and $M\times[-1,0)$ are connected, so each lies in one of $\Omega_\pm$. Every neighbourhood of a point
of $M=\partial\Omega_+=\partial\Omega_-$ meets both, so the halves lie in different ones. Write
$S_\delta^\pm=S_\delta\cap\Omega_\pm$.
\end{itemize}

\emph{Step 2: the two shells.} Let $A_\pm=\{x\in\Omega_\pm:\delta\le d(x)\le r\}$.
\begin{itemize}
\item \emph{Inner side.} Put $O=\Omega_-$. It is bounded and open, and for $x\in O$ the distance to $O^c$ equals
$d(x)$, with the same nearest points (they lie on $\partial O=M$). So the critical points of the distance to $O^c$
in $O$ are exactly the critical points of $d$ in $\Omega_-$. The criticality notions agree:
\cite{ChazalLieutier07} use the vanishing of their gradient field, which holds iff $x$ lies in the convex hull of
its nearest boundary points~\cite[Lemma~2.2]{ChazalLieutier07} (see also~\cite[\S2.2]{CCSL09dcg}).
\item Since $r<\wfs(M)$, the shell $A_-$ contains no critical point. The isotopy lemma~\cite[Prop.~3.4]{ChazalLieutier07}, which they reproduce from
Grove's critical point theory for distance functions, gives a homeomorphism $A_-\cong S^-_\delta\times[\delta,r]$.
\item \emph{Outer side.} Put $O=\Omega_+\cap B(0,\mathcal R)$ with $\mathcal R$ large. For $x\in O$ the distance to
$O^c$ is $\min(d(x),\mathcal R-|x|)$. Take $\mathcal R$ larger than $2r$ plus the largest $|y|$, $y\in M$. Then the
points of $O$ with distance to $O^c$ in $[\delta,r]$ form two disjoint pieces: $A_+$, where $d<\mathcal R-|x|$ strictly, so
the distance equals $d$ with the same nearest points,
and a shell near the sphere $|x|=\mathcal R$, where it equals $\mathcal R-|x|$.
\item Neither piece contains a critical point. In $A_+$ this is because $r<\wfs$. Near the sphere the only nearest
point is the radial projection.
\item So Prop.~3.4 gives a homeomorphism of the union with $R^{-1}(\delta)\times[\delta,r]$, where
$R^{-1}(\delta)=S^+_\delta\sqcup\{|x|=\mathcal R-\delta\}$ has exactly two components ($S^+_\delta\cong M$ is connected).
The components of the product are $C\times[\delta,r]$ for the two components $C$. The shell near the sphere is one
component of the union, so $A_+$ is the other: $A_+\cong C\times[\delta,r]$ with $C$ connected. Step~3 shows that
$C\cong S^+_\delta$.
\item We use Prop.~3.4 in its stated closed-shell form: for a bounded open $O$ and $R=d(\cdot,O^c)$, if the shell
$\{x\in O:\rho_1\le R(x)\le\rho_2\}$ contains no critical point of $R$, it is homeomorphic to
$R^{-1}(\rho_1)\times[\rho_1,\rho_2]$~\cite[Prop.~3.4]{ChazalLieutier07}.
\end{itemize}

\emph{Step 3: gluing.}
\begin{itemize}
\item \emph{The boundary of $A_\pm$.} $A_\pm$ is a compact 3-manifold in $\R^3$, so by invariance of domain its manifold
boundary is its frontier in $\R^3$.
  \begin{itemize}
  \item Points with $\delta<d<r$ are interior points.
  \item A point with $d=\delta$ is a limit of points with $d<\delta$: move along the segment to a nearest point.
  \item A point $x$ with $d=r$ is a limit of points with $d>r$. Since $x$ is not critical, there is a unit vector $u$
  with $\langle u,x-y\rangle>0$ for all $y\in\Gamma(x)$. The one-sided derivative of $d$ in direction $u$ is
  $\min_{y\in\Gamma(x)}\langle u,(x-y)/|x-y|\rangle>0$ (standard formula for the directional derivative of a distance
  function; not separately checked against a source).
  \end{itemize}
  So $\partial A_\pm=S^\pm_\delta\sqcup S^\pm_r$, and the connected closed set $S^\pm_\delta$ is one boundary component.
  Every boundary component of $C\times[\delta,r]$ is a copy of $C$, so $C\cong S^\pm_\delta$.
\item A homeomorphism between manifolds with boundary maps boundary components to boundary components. So the
homeomorphism of Step~2 takes $S^\pm_\delta$, a boundary component of $A_\pm$, onto one end of the product. Composing
with a flip of $[\delta,r]$ if necessary, and with the inverse of its restriction to $S^\pm_\delta$, we get
$A_\pm\cong S^\pm_\delta\times[0,1]$ with $S^\pm_\delta$ going to $S^\pm_\delta\times\{0\}$.
\item Then $\bar N_r(M)=\bar N_\delta(M)\cup A_+\cup A_-$ is $M\times[-1,1]$ with a product collar attached to each
boundary component along that component. That is homeomorphic to $M\times[-2,2]\cong M\times[-1,1]$, with $M$ still
going to $M\times\{0\}$.
\end{itemize}

(a$'$) The interior of $\bar N_r(M)$ is $N_r(M)$: by the computation in Step~3, a point with $d=r$ is a limit of
points with $d>r$.
\begin{itemize}
\item By (a), there is a homeomorphism $g\colon M\times(-1,1)\to N_r(M)$.
\item $g^{-1}(\bar N_d(M))$ is compact, so it lies in $M\times[-a,a]$ for some $a<1$. Take $a<b<1$ and
$K=M\times[-b,b]$, a polyhedral 3-manifold with boundary in the PL 3-manifold $M\times(-1,1)$.
\item The compact set $g(K)$ and the closed set $\R^3\setminus N_r(M)$ are disjoint, and so are the compact sets $g(\partial K)$ and $\bar N_d(M)$. Let
$\varepsilon>0$ be smaller than both distances. By Moise's approximation theorem~\cite[Thm.~35.2]{Moise77} there is a
PL homeomorphism $f$ of $K$ into $\R^3$ with $|f-g|<\varepsilon$ on $K$. Put $U=f(K)$.
\item Then $U$ is a compact PL submanifold, $U\cong_{PL}M\times[-1,1]$, and $U\subseteq N_r(M)=\inter\bar N_r(M)$.
\item $\bar N_d(M)\subseteq\inter U$. Let $x\in\bar N_d(M)$. The straight-line homotopy from $g|_{\partial K}$ to
$f|_{\partial K}$ avoids $x$, so each of the two boundary surfaces has $x$ in its bounded complementary component for $f$
exactly when it does for $g$ (the mod-2 degree of the radial projection from $x$ is a homotopy invariant). For a compact
connected 3-manifold in $\R^3$ with two boundary surfaces, the interior is the set of points of depth one with respect to
them (Definition~\ref{def:solid}; standard, by Alexander duality). Since $x\in g(\inter K)$, its depth is one for $g$,
hence for $f$, so $x\in\inter U$.
\item $f$ need not send $M$ to a level, which is why the collar condition of Section~\ref{sec:uniqueness} does not ask
for that.
\end{itemize}

(b) The zones lie in $N_d(M)\subseteq\inter U$, and $M\subseteq N_d(M)\subseteq\inter U$. The $U_i$ are pairwise
disjoint because they lie in the disjoint sets $\bar N_r(M_i)$. It remains to show that the inclusion
$M\hookrightarrow U$ is a homotopy equivalence.
\begin{itemize}
\item By (a), the composite $M\hookrightarrow U\hookrightarrow\bar N_r(M)$ is a homotopy equivalence. So
$\pi_1(U)\to\pi_1(\bar N_r(M))$ is onto.
\item Both groups are isomorphic to $\pi_1(M)$, a surface group of genus at least one. Surface groups are Hopfian,
so a surjection from one onto an isomorphic copy is an isomorphism. Hence $\pi_1(M)\to\pi_1(U)$ is an isomorphism.
\item $M$ and $U\simeq M$ are aspherical, so by Whitehead's theorem the inclusion is a homotopy equivalence.
\end{itemize}
This uses (a) as well as (a$'$), and it uses the genus hypothesis: for a sphere the argument fails.
Standard facts used without separate checking against a source: the Jordan--Brouwer separation theorem,
Tarski--Seidenberg, regular-neighbourhood theory (the regular neighbourhood of a two-sided closed surface is a product),
invariance of boundary for topological manifolds, attaching a product collar along a boundary component, that surface
groups are Hopfian (they are residually finite), asphericity of closed surfaces of positive genus, Whitehead's
theorem, homotopy invariance of the mod-2 degree and its use to decide on which side of an embedded closed surface a
point lies, and Alexander duality for the components of the complement of two disjoint closed surfaces.
\end{proof}

\subsection{An explicit lower bound for polyhedral surfaces}
For $a\in M$, the \emph{conical radius} $\rho(a)$ is the supremum of the $r>0$ such that $M\cap B(a,r)=C\cap
B(a,r)$ for some closed cone $C$ with apex $a$.

\begin{lemma}[Cone lemma]\label{lem:cone}
Let $a\in M$, $r<\rho(a)$, and $p\notin M$ with $|p-a|<r/2$. Then every nearest point $x\in\Gamma(p)$ satisfies
$\langle p-x,\,p-a\rangle=|p-x|^2$. In particular $p$ is not critical.
\end{lemma}
\begin{proof}
\begin{itemize}
\item \emph{Nearest points lie in the ball.} $d(p)\le|p-a|<r/2$, so every $x\in\Gamma(p)$ has $|x-a|<r$.
\item \emph{The ray through $x$ stays in $M$.} The points $y_s=a+s(x-a)$, for $s$ near $1$, lie in the cone and
in $B(a,r)$, hence in $M$. So $s\mapsto|p-y_s|^2$ has a local minimum at $s=1$, which gives $\langle
p-x,x-a\rangle=0$.
\item \emph{Conclusion.} $\langle p-x,p-a\rangle=|p-x|^2+\langle p-x,x-a\rangle=d(p)^2>0$ for all
$x\in\Gamma(p)$. So $0$ is not in the convex hull of the vectors $p-x$.
\end{itemize}
\end{proof}

\begin{corollary}\label{cor:lambda}
Let $\lambda(M)=\inf_{x\in M}\sup_{a\in M}\bigl(\rho(a)/2-|x-a|\bigr)$. Then $d$ has no critical point $p$ with
$0<d(p)<\lambda(M)$, so $\wfs(M)\ge\lambda(M)$. If $M$ is a finite union of triangles, $\lambda(M)>0$.
\end{corollary}
\begin{proof}
\begin{itemize}
\item \emph{The bound.} Given $p$ with $d(p)<\lambda$ and $x\in\Gamma(p)$, pick $a$ with
$\rho(a)/2-|x-a|>d(p)$. Then $|p-a|<\rho(a)/2$, and Lemma~\ref{lem:cone} applies.
\item \emph{Positivity.} For a polyhedron, $\rho(a)>0$ for every $a$, by the local conical structure of
polyhedra (standard; not separately checked against a source). The function $g(x)=\sup_a(\rho(a)/2-|x-a|)$
is $1$-Lipschitz and finite, since $\rho\le\operatorname{diam}M$, and $g(x)\ge\rho(x)/2>0$. $M$ is compact, so
$\min g>0$.
\end{itemize}
\end{proof}

Lemma~\ref{lem:cone} and the pointwise form of Corollary~\ref{cor:lambda} are machine-checked in Lean~4 with
Mathlib, for subsets of an arbitrary real inner product space (Appendix~\ref{sec:lean}).

\subsection{A certified lower bound for $\wfs$}
For a triangulated reference surface $M_{PL}$ we compute $L$ with ``no critical value in $(c\cdot\mathrm{tol},L)$''
by a best-first octree. We call $L$ \emph{certified} in this qualified sense only: the computation is in double
precision; containment of a point in a feature is decided with a tolerance $\mathrm{tol}=10^{-10}$ times the model
size; and the search is restricted to distances below a cut-off $R$.
\begin{itemize}
\item \emph{Boxes near $M$} are certified by Lemma~\ref{lem:cone}, using a lower bound for $\rho(a)$ at a
candidate point $a$. The bound comes from the distance from $a$ to the nearest feature (planar region, boundary
segment, vertex) that does not contain $a$. It is sound for planar regions whose directed boundary segments cover
their boundary, after a global coplanarity re-check of each merged region.
\item \emph{Boxes farther away} are certified by a direction test: a small linear program finds a unit vector $u$
with $\langle u,p-x\rangle>0$ for all candidate nearest points $x$ of all points $p$ of the box. The candidates
are the projections of the box centre onto the relevant features.
\item \emph{Stopping rule.} The search stops at the first box that is too small to split further and cannot be
certified. Its distance from $M$ is the certified bound $L$.
\end{itemize}
Table~\ref{tab:wfs} lists the results. Lemma~\ref{lem:collar} needs no critical value in $(0,r]$. The computed
bound covers $(c\cdot\mathrm{tol},L)$; the interval $(0,c\cdot\mathrm{tol}]$ is covered by
Corollary~\ref{cor:lambda} whenever $\lambda(M)>c\cdot\mathrm{tol}$, which we have not checked separately.

\begin{table}[h]
\centering
\begin{tabular}{lcc}
\toprule
Part & exact $\wfs$ (smooth part) & certified lower bound for $M_{PL}$ \\
\midrule
cube, side 10 & 5 & 4.530 \\
plate $4\times4\times1$ & 0.5 & 0.436 \\
plate $4\times4\times0.3$ & 0.15 & 0.137 \\
L-shape (wall 2) & 1 & 0.894 \\
L-shape, inner fillet 0.2 & 1 & 0.894 \\
block with slot of width 0.2 & 0.1 & 0.086 \\
cylinder $R=3$, $H=10$ & 3 & 2.512 \\
\bottomrule
\end{tabular}
\caption{Lower bounds for $\wfs$ of the triangulated parts, certified in the qualified sense above. The exact
column is the $\wfs$ of the smooth part; for a triangulated cylinder $\wfs(M_{PL})$ is slightly below $3$. The inner
fillet caps the reach at $0.2$ but leaves the bound unchanged.}\label{tab:wfs}
\end{table}

\section{Existence}\label{sec:existence}

We give two constructions that turn hypotheses on the data into a realization. Neither assumes that a perturbed
surface embeds as a whole. Theorem~\ref{thm:lemma0} does assume (P0) that each original patch is embedded, and two
of its hypotheses are non-local. Propositions~\ref{prop:B} and~\ref{prop:Bprime} assume that the re-trimmed model is
embedded, which is checked computationally on the model itself.

\subsection{Fans and cones}\label{sec:fans}

\subsubsection*{Setting}
$K$ is regular: every face is a disk bounded by a single wire, and every edge has two distinct vertices. The theorem
and its proof (Appendix~\ref{sec:proof51}) cover only this case. For B-rep complexes, natural readings of the
hypotheses exist: at a seam edge E3$'$ refers to the two sheets of the face; at a closed edge V1 counts edge ends; for
a face with inner wires, P0 asks that $Q_f$ be homeomorphic to the face. But the proof does not cover them. For
example, at a seam the two fan sheets of one face glue to an annulus. These cases are not proved here.
$\mathrm{dev}^*_e(t)$ is the deviation of the boundary points $b_{f,e}(t)=\partial P_f\cap\Pi(t)$ from $\gamma_e(t)$ in
the normal plane $\Pi(t)$; it is at most $\mathrm{dev}_e$.

\subsubsection*{Auxiliary sets}
\begin{itemize}
\item \emph{Normal disks.} For an edge $e$, choose parameters $0<a_e<b_e<1$ and a $C^1$ radius function $\rho_e(t)$.
$\rho_e$ is constant on the stubs $[t_v,a_e]$, $[b_e,t_w]$ and on a neighbourhood of $a_e$ and $b_e$; here $t_v$ is
the last parameter with $\gamma_e(t)\in B(p_v,r_v)$. $D_e(t)$ is the closed disk of radius $\rho_e(t)$ in $\Pi(t)$
centred at $\gamma_e(t)$, and $T_e=\bigcup_{t\in[a_e,b_e]}D_e(t)$.
\item \emph{Vertex regions.} For a vertex $v$, choose $r_v$. The vertex region $W_v$ is the union of the ball
$B(p_v,r_v)$ and, for each edge $e$ at $v$, the disks $D_e(t)$ for $t$ between $t_v$ and $a_e$ (or $b_e$).
\end{itemize}

\subsubsection*{Hypotheses}
E1--E3$'$ are required on $[t_v,t_w]$, the stubs included.
\begin{description}
\item[E1] $\mathrm{dev}^*_e(t)<\rho_e(t)<\min(\tau_e,\operatorname{reach}\gamma_e|_{[a_e,b_e]})$.
\item[E2] For every face $f>e$:
  \begin{enumerate}
  \item $\psi_f<\pi/2$, where $\psi_f$ is the largest angle between $\gamma_e'$ and $T P_f$ inside the tube;
  and $\rho_e k_f<1$, where $k_f=\kappa_f/\cos\psi_f$;
  \item $t\mapsto\pi_e(\beta_{f,e}(t))$ is strictly monotone. A sufficient condition is $L_e+\kappa_\gamma
  \mathrm{dev}^*_e<1$, where $L_e$ is the arc-length Lipschitz constant of $\beta_{f,e}-\gamma_e$;
  \item \emph{no return:} $\dist(\gamma_e(t),P_f\setminus C_f)>\rho_e(t)$. Here $C_f$ is the union of the section
  arcs of $P_f$ from $b_{f,e}$ to their exit points, of length at most
  $\ell_f=(\mathrm{dev}^*+\sqrt{\mathrm{dev}^{*2}+\mu\rho^2})/\mu$ with $\mu=1-\rho k_f$.
  \end{enumerate}
\item[E3$'$] The exit points $k_{f,e}(t)$ and $k_{g,e}(t)$ of the two face sections on $\partial D_e(t)$ are distinct.
A sufficient condition (E3) is $\alpha_e(t)>\varepsilon_f+\varepsilon_g$, with
$\varepsilon=\arcsin\min\{1,(\mathrm{dev}^*_e+k\ell^2/2)/\rho_e\}$. Here $\alpha_e$ is the angle between the inward
tangents of the two sections.
\item[V0] $\gamma_e$'s end at $v$ and the corners of the patches at $v$ lie in $\inter B(p_v,r_v)$.
\item[V1] (a) $W_v\subseteq\inter\Z(v)$, and $W_v$ is \emph{strictly star-shaped} with respect to $p_v$: for every
$x\in\partial W_v$, the vector $p_v-x$ lies in the interior of the tangent cone of $W_v$ at $x$, and each ray from
$p_v$ meets $\partial W_v$ exactly once. (The second property does not follow from the first at junction points, where a
ray could graze $\partial W_v$, so it is stated separately.)
(b) \emph{(non-local)} The construction outside $W_v$ meets $\partial W_v$ in one simple closed curve $\lambda_v$.
\begin{itemize}
\item Each incident edge contributes one point and each incident face one arc, in the cyclic order of the link of $v$.
\item The outside pieces cross $\partial W_v$ transversally: along the smooth strata of $\partial W_v$ (the ball and the
stub cylinders), and along the junction curves where a stub cylinder meets the ball sphere.
\item $\lambda_v$ meets these junction curves transversally and avoids the corners of $\partial W_v$.
\end{itemize}
\item[S1] (a) The $W_v$ are pairwise disjoint. (b) The $T_e$ are pairwise disjoint. (c) $\gamma_e$ and $T_e$ miss $W_w$
for $w\notin\partial e$. (d) $P_f$ misses $T_e$ for $e\not<f$, and misses $W_v$ for $v\not<f$. (e) Outside
$\bigcup T\cup\bigcup W$, the patch pieces are pairwise disjoint. (f) $T_e\cap W_v=D_e(a_e)$ for $v<e$, and each
piece outside $W_v$ meets $W_v$ only in $\lambda_v$.
\item[P0] \emph{(non-local)} Each $P_f$ is embedded, and $P_f\cap(\bigcup T\cup\bigcup W)$ is a collar of $\partial
P_f$, so that $Q_f=P_f\setminus\inter(\bigcup T\cup\bigcup W)$ is a closed disk.
\end{description}
V1(b) and P0 are not consequences of local data. In the models of Section~\ref{sec:experiments} they were verified
only for the gapped cube, implicitly, through an exact check that the constructed surface is embedded. The
sampled checker does not test V1(b), P0, E2.3 beyond S1(d), or S1(e) for adjacent faces.

\begin{theorem}[Existence by fans and cones; conditional]\label{thm:lemma0}
Under E1--E3$'$, V0, V1, S1 and P0, the tolerant B-rep has a realization.
\end{theorem}
The proof is in Appendix~\ref{sec:proof51}. In outline:
\begin{itemize}
\item $h$ follows $\gamma_e$ along the edges and puts fans of radii in the normal disks.
\item The patches, with their collars removed, are kept away from the tubes.
\item Cones over the curves $\lambda_v$ fill the vertex regions.
\item Injectivity comes from the disjoint normal disks, E3$'$, star-shapedness and S1/P0.
\item The resulting map is bi-Lipschitz, so a fine secant approximation is a PL realization.
\end{itemize}

\begin{example}
On a cube whose faces are pulled apart by gaps, the construction was carried out exactly in rational
arithmetic: 216 triangles, with zone membership and embeddedness checked exactly
(\texttt{compute/lemma0\_box.py}).
\end{example}

\subsection{Re-intersection}

\begin{proposition}[Planar, degree three]\label{prop:B}
Suppose every face is a disk bounded by a single wire and is planar, $P_f\subseteq\Pi_f$, and every vertex lies on exactly three faces whose planes are
independent. Let $c_v$ be the common point of the three planes, $\ell_e=[c_v,c_w]$ for an edge $e=vw$, and
$\tilde P_f$ the polygon on the $c$'s of the vertices of $f$, in the order of $\partial f$. Assume:
\begin{itemize}
\item $c_v\in\Z(v)$, $\ell_e\subseteq\Z(e)$ and $\tilde P_f\subseteq\Z(f)$;
\item the $\tilde P_f$ are simple polygons that meet exactly in their common edges and vertices.
\end{itemize}
Then the cellular map onto $\bigcup\tilde P_f$ is a realization.
\end{proposition}
\begin{proof}
It is a PL homeomorphism onto $\bigcup\tilde P_f$, and each cell lands in its zone.
\end{proof}

\begin{proposition}[Curved re-intersection]\label{prop:Bprime}
Suppose there is a \emph{re-trimmed model}: compact patches $\tilde P_f$, curves $\tilde\gamma_e$ and points $c_v$, with
$\Sigma=\bigcup_f\tilde P_f$, satisfying the following.
\begin{description}
\item[P] Each face has a $C^2$ regular parametrization $X_f\colon U_f\to\R^3$, where $U_f$ is open in $\R^2$ or in a
quotient $\R^2/\Lambda_f$ by a lattice of periods (this is how seams are represented). There is a compact region
$D_f\subset U_f$ with $X_f(D_f)=\tilde P_f$ and $X_f$ injective on $D_f$. Singular parametrizations, such as sphere
poles and cone apices, are excluded.
\item[C] There is a homeomorphism $\phi\colon|K|\to\Sigma$ mapping each open cell of $K$ onto the corresponding open cell
of $\Sigma$. The open cells of $\Sigma$ are the points $c_v$, the curves $\tilde\gamma_e$ minus their end points, and the
patches minus all their edge curves (seams included).
\item[R] Let $G_f\subset D_f$ be the \emph{edge graph} of $f$: the preimage under $X_f$ of all edge curves of $f$,
including seams, which in the quotient parametrization lie inside $D_f$ rather than on its boundary. $G_f$ is a finite
graph. Its vertices are the preimages of the points $c_v$. Its edges are simple $C^2$ regular arcs, or loops at a
vertex, meeting only at vertices. At each vertex of $G_f$, every sector of $D_f$ between two consecutive arcs has angle
in $(0,2\pi)$. A smooth passage, with angle $\pi$, is allowed: for example a seam vertex on a smooth boundary curve.
\item[B1--B3] $c_v\in\Z(v)$, $\tilde\gamma_e\subseteq\Z(e)$ and $\tilde P_f\subseteq\Z(f)$.
\item[T] At every interior point of every edge curve $\tilde\gamma_e$, the inward tangent half-planes of the two adjacent
patches (or of the two sides of a seam) are distinct: the surface does not fold back on itself. At every vertex $c_v$,
the tangent cone of $\Sigma$ is an embedded disk. That cone is the union of the tangent sectors $dX_f(\text{sector})$,
taken over all sectors of the edge graphs $G_f$ at preimages of $c_v$. The condition means these sectors meet only along
common boundary rays, in the cyclic order of the link of $v$. (Together with the $C^1$ regularity up to the vertices,
this also controls the edge condition at the end points of the edges.)
\end{description}
Then the tolerant B-rep has a realization.
\end{proposition}

\begin{remark}
\begin{itemize}
\item By C, $\Sigma$ is a disjoint union of closed surfaces, so it bounds a solid in the sense of
Definition~\ref{def:solid}.
\item C covers seams, closed edges and faces with inner wires.
\item The edge part of T asks only that the surface does not fold. Transversal faces satisfy it, and so do tangent (G$^1$)
edges, where the two half-planes are opposite. So the proposition applies in principle to filleted models, provided a
re-trimmed model with the stated properties is available. Our computation builds the re-trimmed model by
re-intersecting extended surfaces, which fails for tangent faces (Section~\ref{sec:experiments}); that is a limitation
of the computation, not of the proposition.
\end{itemize}
\end{remark}
The proof is in Appendix~\ref{sec:proofBprime}.

\begin{remark}[Checking the hypotheses]
\begin{itemize}
\item The re-trimmed model can be computed by extending the face surfaces and rebuilding the solid
(OpenCascade's \texttt{BOPAlgo\_MakerVolume}). The hypotheses are then checked on samples
(\texttt{compute/reintersect\_curved.py}), except the vertex part of hypothesis T, which is not
checked. That matters for the cylinder and the pipe, whose seam vertices lie on only two surfaces. The PL step of the
proof is not carried out numerically; it is an existence argument.
\item For the planar case, the hypotheses of Proposition~\ref{prop:B} for the sewn box and the sewn plate were also
re-checked in exact rational arithmetic, including embeddedness.
\item Results are in Section~\ref{sec:experiments}.
\end{itemize}
\end{remark}

\section{An obstruction}\label{sec:obstruction}

\subsection{Setting}
\begin{itemize}
\item $f\ne g$ are faces whose closed cells are disks bounded by single wires.
\item Their data are planar: $P_f\subseteq\Pi_f$ and $P_g\subseteq\Pi_g$, with the planes not parallel.
\item $\lambda=\Pi_f\cap\Pi_g$, and $\alpha\in(0,\pi/2]$ is the angle between the planes.
\item $n_f,n_g$ are unit normals, and $m_f$ is a unit vector in $\Pi_f$ orthogonal to $\lambda$ (likewise $m_g$).
\item $\pi_f$ is the orthogonal projection onto $\Pi_f$, and $U_f=\bigcup_{e<f}\Z(e)$, which contains the zones
of the vertices of $f$.
\item $\ell_f$ is the boundary loop of $f$: the arcs $\beta_{f,e}$ in the cyclic order of $\partial f$, with
consecutive arcs joined by the segment between their end points.
\end{itemize}

\begin{theorem}[Obstruction]\label{thm:obstruction}
Suppose there are $X\in\lambda$ and $\rho>0$ such that, with $J_f=X+[-\rho,\rho]m_f$ and $J_g=X+[-\rho,\rho]m_g$:
\begin{description}
\item[O0] $\beta_{f,e}\subseteq\Z(e)$ for every edge $e$ of $f$, and at every vertex $v$ of $f$ the end points of
the two arcs of $f$ at $v$ lie in $\Z(v)$; likewise for $g$.
\item[O1] $J_f\cap\ell_f=\emptyset$, and the winding number of $\ell_f$ in $\Pi_f$ about the points of $J_f$ is
odd; likewise for $g$.
\item[O2] $J_f\cap\cl\,\pi_f(U_f)=\emptyset$ and $J_g\cap\cl\,\pi_g(U_g)=\emptyset$.
\item[O3] $\Z(e)$ is simply connected for every edge $e$ of $f$ or $g$.
\item[O4] $\rho\sin\alpha>\tau_f\cos\alpha+\tau_g$ and $\rho\sin\alpha>\tau_g\cos\alpha+\tau_f$.
\end{description}
Then the tolerant B-rep has no realization.
\end{theorem}

\begin{proof}
Suppose $h$ is a realization, and let $\varphi=\pi_f\circ h|_{\bar f}$.

\emph{Step 1: over $J_f$, $h(\bar f)$ lies in the slab.}
\begin{itemize}
\item Let $y\in J_f$. We have $h(\bar f)\subseteq\Z(f)=N_{\tau_f}(P_f)\cup U_f$, and by O2 no point of $U_f$ lies over
$y$.
\item Since $P_f\subseteq\Pi_f$, $\pi_f^{-1}(y)\cap N_{\tau_f}(P_f)\subseteq\sigma_y:=y+(-\tau_f,\tau_f)n_f$.
\item So $h(\bar f)\cap\pi_f^{-1}(y)=h(\bar f)\cap\sigma_y$. These points are images of interior points of $f$,
because $h(\partial f)\subseteq U_f$.
\end{itemize}

\emph{Step 2: degree one.}
\begin{itemize}
\item Parametrize $\ell_f$ over $\partial f$: a small arc $A_v$ around each vertex goes to the joining segment, and
the rest of each edge goes to $\beta_{f,e}$. Shrink $A_v$ until $h(A_v)\subseteq\Z(v)$.
\item Build a homotopy in $U_f$ from $h|_{\partial f}$ to this parametrization:
  \begin{itemize}
  \item on $A_v$, the straight-line homotopy inside the convex ball $\Z(v)$;
  \item on the rest of an edge $e$, the boundary of the square is mapped into $\Z(e)$ (by O0, monotonicity, and
  $h(\bar e)\subseteq\Z(e)$), so it extends over the square by O3.
  \end{itemize}
\item By O2 the projected homotopy avoids every $y\in J_f$. So $\varphi|_{\partial f}$ and $\ell_f$ have the same
winding number about $y$, which is odd by O1. Hence the mod-2 degree of $\varphi$ at $y$ is $1$.
\item Fix a triangulation of $\bar f$ on which $h$ is linear. For $y\notin\varphi(\text{1-skeleton})$, which
excludes only finitely many points of $J_f$ once the slice in Step 3 is generic, $\#(h(\bar f)\cap\sigma_y)$ is
odd.
\end{itemize}

\emph{Step 3: a crossing arc.}
\begin{itemize}
\item Let $\Sigma$ be the plane through $X$ orthogonal to $\lambda$; it contains $n_f,m_f,n_g,m_g$.
\item Move $X$ along $\lambda$ and shrink $\rho$ slightly to generic values; O1, O2 and O4 are open conditions.
Genericity means: $\Sigma$ contains no vertex of the triangulations and meets their edges transversally, and the
lines $\langle x-X,m_f\rangle=\pm\rho$, $\langle x-X,m_g\rangle=\pm\rho$ contain no vertex.
\item In $\Sigma$, the rectangle $R_f=J_f+[-\tau_f,\tau_f]n_f$ meets $h(\bar f)$ in a compact 1-manifold $M_f$.
\item $M_f$ misses $h(\partial f)$, since $\pi_f(R_f)=J_f$ and O2 holds. It misses the long sides of $R_f$, by
Step 1. So its boundary lies on the two short sides.
\item A generic fibre $\sigma_y$ meets $M_f$ an odd number of times (Step 2). Circle components, and arcs with both
ends on one short side, meet it an even number of times.
\item Hence some arc $a_f\subseteq M_f$ joins the two short sides of $R_f$. Likewise there is an arc $a_g$ for $g$.
\end{itemize}

\emph{Step 4: the arcs meet.}
\begin{itemize}
\item In $\Sigma$, use the coordinates $u=\langle x-X,n_f\rangle$ and $w=\langle x-X,n_g\rangle$.
\item On the short sides of $R_f$, $|w|\ge\rho\sin\alpha-\tau_f\cos\alpha>\tau_g$ by O4, with opposite signs on
the two sides. So $a_f$ contains a sub-arc in the parallelogram $Q=\{|u|\le\tau_f,|w|\le\tau_g\}$ joining
$w=-\tau_g$ to $w=\tau_g$. Symmetrically, $a_g$ contains a sub-arc in $Q$ joining $u=-\tau_f$ to $u=\tau_f$.
\item By O4, $Q$ lies in both rectangles.
\item Two paths in a square joining the two pairs of opposite sides intersect (a standard consequence of the Brouwer
fixed point theorem; not separately checked against a source). So $h(\bar f)\cap h(\bar g)\ne\emptyset$ at a point that is
the image of interior points of both $f$ and $g$ (Step 1). This contradicts injectivity.
\end{itemize}
\end{proof}

\begin{remark}
\begin{itemize}
\item \emph{Which hypotheses bind.} Face tolerances are tiny (OpenCascade's default is $10^{-7}$), so O4 is mild.
The binding conditions are O1 and O2: the patches must cross away from the shadows of the edge and vertex zones.
\item \emph{Curved patches.} The proof uses planarity in Steps 1 and 2. A version for $C^2$ patches with $\tau_f$
below their reach should follow with tubular coordinates; it is not written.
\item \emph{Standard facts used without separate checking against a source:} mod-2 degree of PL maps and its relation
to the winding number, general position for PL maps and generic planes, parity of crossings of a segment with a PL
circle, homotopy invariance of the winding number, the nerve lemma (for the O3 test in the computations), and the
square lemma cited above.
\end{itemize}
\end{remark}

\section{Experiments with OpenCascade}\label{sec:experiments}

\subsection{Set-up}
\begin{itemize}
\item \emph{Software.} All experiments use OpenCascade 8.0.1 through \texttt{cadquery-ocp}. Scripts and results are
in the repository~\cite{repo}.
\item \emph{The two checks.} ``BRepCheck'' is \texttt{BRepCheck\_Analyzer(shape).IsValid()}. ``BOP check'' is
\texttt{BRepAlgoAPI\_Check(shape, true, true)}, which also tests geometry and self-interference.
\item \emph{Sewn models} imitate imported data. Every face of an exact part is moved by a random vector of length
between $g/2$ and $g$, and the faces are sewn back with \texttt{BRepBuilderAPI\_Sewing} at tolerance $4g$. Sewing
sets each edge tolerance to $1.05$ times its own deviation estimate, which is $1.048$--$1.05$ times our measured
deviation.
\item \emph{Scope.} These are a few parts in one kernel version, not a survey of industrial models.
\end{itemize}

\subsection{The checks do not decide existence or uniqueness}
Table~\ref{tab:occ} summarises.

\begin{table}[h]
\centering\small
\begin{tabular}{>{\raggedright}p{4.4cm}cc>{\raggedright\arraybackslash}p{5.0cm}}
\toprule
Model & BRepCheck & BOP check & What is true \\
\midrule
twisted triangular prism, two faces with linked boundaries ($\tau=0.081$) & valid & self-intersection & no
realization (linking number $-1$) \\
thick pipe, tolerance $1.3\times$ hole radius & valid & valid & two realizations with non-homeomorphic solids
(knot determinant $3$ vs.\ $1$) \\
sewn wedges, $30^\circ$ and $10^\circ$, $g=0.01$ & valid & valid & no realization
(Theorem~\ref{thm:obstruction}) \\
sewn plate, thickness $0.05$, $g=0.01$ & valid & invalid & realizable (Proposition~\ref{prop:B}) \\
filleted L-shape, $\tau=0.1$ & valid & self-intersection & $\tau+\mathrm{defl}<$ certified $\wfs$ bound of
$M_{PL}$ \\
long cylinder, $\tau=3.3$ & valid & valid & $\tau\ge\wfs$(nominal) (critical point exhibited) \\
\bottomrule
\end{tabular}
\caption{OpenCascade's checks against what the theory shows.}\label{tab:occ}
\end{table}

\subsubsection*{Twisted prism}
\begin{itemize}
\item Two non-adjacent triangular faces $A,B$ have boundary loops with linking number $-1$ (by crossing count; a
numerical Gauss integral agrees). The loops are $0.4037>2\tau=0.1615$ apart, so the unions $T_A$, $T_B$ of
their edge and vertex zones are disjoint.
\item In a realization, each $h(e)$ lies in the convex capsule $C_e=N_\tau(e)$, which contains the end balls (all
tolerances are equal here). Homotope $h(e)$ rel end points to a segment inside $C_e$, and the vertices to the
nominal ones inside the balls. This gives a homotopy of $h(\partial A)$ to $\partial A$ inside $T_A$, so in the
complement of $T_B$; likewise for $B$.
\item The linking number is invariant under such homotopies, so $\mathrm{lk}(h(\partial A),h(\partial B))=-1$, and
$h(\partial A)$ must meet the disk $h(\bar B)$. That contradicts embeddedness.
\end{itemize}

\subsubsection*{Knotted pipe}
\begin{itemize}
\item One realization keeps the nominal geometry. The other replaces the inner face by a thin tube, of radius
$0.005$, around a centreline that follows a long trefoil.
\item The script checks the zones (maximum distance $0.983\le1.3$), embeddedness ($\varepsilon\kappa_{\max}=0.31$;
non-local distance $0.036>2\varepsilon$ outside a window of $0.04<\pi/\kappa_{\max}=0.0504$) and the knot
(determinant $3$).
\item The second solid is the exterior of a knotted arc in a ball, whose
fundamental group is non-abelian: a knot of determinant $3$ admits a non-trivial Fox $3$-colouring, and a
non-trivial $3$-colouring of a diagram defines, through the Wirtinger presentation, a surjection of the knot group onto
the non-abelian group $S_3$ (standard; not separately checked against a source). The nominal
solid is a solid torus. So the two solids are not homeomorphic.
\item This shows that the collar hypothesis of Theorem~\ref{thm:uniqueness} cannot be dropped. Here
$\wfs\le1<1.3$.
\end{itemize}

\subsubsection*{Sewn wedges}
\begin{itemize}
\item At the knife edge of the sewn wedges, the two planar patches cross along a line that lies outside the edge
tube, by $49\%$ ($30^\circ$) and about $3\times$ ($10^\circ$) of $\tau_e$.
\item The hypotheses of Theorem~\ref{thm:obstruction} hold. For the $30^\circ$ wedge the largest feasible $\rho$ at
the positions searched is $1.6\cdot10^{-3}$, while O4 needs only $3.7\cdot10^{-7}$; for the $10^\circ$ wedge it is
$1.6\cdot10^{-2}$.
\item An independent check in exact rational arithmetic confirmed all of O0--O4 for both wedges.
\item On the sewn box and plate and on the exact box and wedge, the checker finds no crossing witness, consistent
with their being realizable.
\end{itemize}

\subsection{Which construction applies}

\begin{table}[h]
\centering\small
\begin{tabular}{>{\raggedright}p{4.0cm}>{\raggedright}p{2.8cm}>{\raggedright}p{2.6cm}>{\raggedright\arraybackslash}p{4.0cm}}
\toprule
Model & fans (Thm.~\ref{thm:lemma0}) & re-intersection & margins of re-intersection (vertex / edge / face) \\
\midrule
exact box, cylinder, pipe, slot, filleted box, $10^\circ$ wedge & pass & pass where applicable & --- \\
L-shape, all edges filleted 0.5 & pass with variable radius$^*$ & not applicable & --- \\
L-shape fused, inner fillet 0.2 & fails (2 cusp vertices, 9 V1(a) failures); passes after merging coplanar
faces & not applicable & --- \\
sewn box $g=0.01$ & fails at vertices & pass & 45\% / 13\% / 4.9\% \\
sewn plate $g=0.01$ & fails at vertices & pass & 45\% / 13\% / 5.9\% \\
sewn cylinder $g=0.01$ & fails at vertices & pass (curved) & 55\% / 35\% / 5.0\% \\
sewn cylinder $g=0.05$ & fails at vertices & pass (curved) & 55\% / 11\% / 8.9\% \\
sewn pipe $g=0.01$ & fails at vertices & pass (curved) & 53\% / 16\% / 5.1\% \\
sewn wedges & fails at vertices & fails (edge outside tube) & no realization (Thm.~\ref{thm:obstruction}) \\
sewn filleted box $g=0.001$ & fails at vertices & computation not applicable (tangent faces) & open \\
\bottomrule
\end{tabular}
\caption{Constructions on exact and sewn models (sampled checks, double precision). $^*$Provisional: the checker
does not yet implement the constant radius near the vertices. Face margins are the smaller of the sampled checker's
value and an independent dense re-check with exact projections (\texttt{results/reintersect\_B3\_dense.txt}); for the
sewn plate only the dense re-check is available.}\label{tab:constr}
\end{table}

\begin{itemize}
\item \emph{Why the fans fail on sewn data.} The fan construction fails at the vertices of sewn models. Sewing sets
$\tau_e=1.05\,\mathrm{dev}_e$, and the tube radius must exceed the deviation; but at a corner of a box three tubes of
radius greater than $\mathrm{dev}$ do not fit disjointly into a vertex ball of radius about $\tau_e$.
\item \emph{Re-intersection.} Re-intersection avoids this. Our computation of the re-trimmed model, which
re-intersects the extended surfaces, needs transversal faces. ``Pass (curved)'' refers to the edge part of hypothesis T
only; the vertex part is not checked.
\item \emph{Sewn models with tangent edges remain open.} On the sewn filleted box, OpenCascade sets each vertex
tolerance equal to the largest incident edge tolerance, the patch corners lie up to $0.85\,\tau_v$ from the vertex
point, and none of our vertex constructions fits.
\end{itemize}

\subsection{Other observations}
\begin{itemize}
\item \emph{Thickness heuristics.} Half the minimum distance between non-adjacent faces is a common estimate of the
safe tolerance. It overestimates the cylinder (5 against $\wfs=3$), and underestimates an L-shape with a $0.05$
inner fillet by a factor of about $97$.
\item \emph{Offsets.} In its default mode, \texttt{BRepOffsetAPI\_MakeOffsetShape} returned solids with wrong
volume (errors from $44\%$ to over a thousandfold, with six shells instead of one in every such case) for a filleted box at
inward offsets larger than the fillet radius but within the $\wfs$-safe range, and BRepCheck accepted them. A plausible
mechanism is in the source of OpenCascade~8.0.1: when the offset radius of a cylindrical, spherical or toroidal fillet
face becomes negative, its offset
surface gets the status \texttt{BRepOffset\_Reversed}, and \texttt{BRepOffset\_MakeOffset::BuildOffsetByArc} skips
such faces (source comment: ``Degenerated or returned faces are not stored''). We have not traced the rest of the
algorithm. These failures appear not to involve tolerances; they look like failures of face-by-face offsetting when faces must
disappear.
\item \emph{Why BRepCheck ignores the size of tolerances.} In OpenCascade~8.0.1 the documentation comment of
\texttt{BRepCheck\_Analyzer} lists the status \texttt{BRepCheck\_InvalidToleranceValue} as ``NYI'' (not yet
implemented) for vertices, edges and faces. In \texttt{BRepCheck\_Analyzer.cxx} the flag that would set this status on a
face is initialised to false and never assigned, and no other file sets the status. So BRepCheck never reports a
tolerance as too large. (This was first noticed in the author's companion book project; we checked the source
ourselves.)
\item \emph{Public documentation.} We did not find a check that tolerances leave the topology unchanged in the
public documentation of OpenCascade, or in third-party documentation of ACIS and Parasolid. That says only that
the documentation does not mention one.
\end{itemize}

\section{Limitations and open problems}\label{sec:open}

\begin{enumerate}
\item \textbf{Citations behind the proofs.} The results quoted from the literature in the proofs have been checked
against the originals. What remains are the standard facts listed, each time, as ``not separately checked against a
source''.
\item \textbf{B-rep complexes in the fan construction.} Theorem~\ref{thm:lemma0} is proved for regular complexes only;
seams, closed edges and faces with inner wires need separate arguments. (Proposition~\ref{prop:Bprime} covers them.)
\item \textbf{Spheres.} Theorem~\ref{thm:uniqueness} excludes shells of genus $0$. Waldhausen's horizontality
argument does not apply to them. The PL Schoenflies theorem~\cite[Thm.~17.12]{Moise77} is the natural tool for the sphere
case; the argument is not written.
\item \textbf{Embeddings.} We prove that the images of realizations are ambient isotopic, not that the embeddings
are isotopic as maps. The two differ by a self-homeomorphism of the surface.
\item \textbf{Tight vertices.} We do not know whether sewn models with tangent edges and OpenCascade's vertex
tolerances (equal to the largest incident edge tolerance) have realizations. Two directions are open: a vertex
construction that uses the edge tubes as well as the ball, or a vertex version of the obstruction.
\item \textbf{Re-trimmed models for tangent edges.} Proposition~\ref{prop:Bprime} allows tangent (G$^1$) edges and mixed
vertices, but our computation of a re-trimmed model re-intersects extended surfaces and fails when faces are tangent. A
construction of re-trimmed models for such data is open.
\item \textbf{Curved obstruction.} Theorem~\ref{thm:obstruction} is stated for planar patches.
\item \textbf{Rigour of the computations.}
  \begin{itemize}
  \item The $\wfs$ bounds and the hypothesis checks use double precision and sampling. Exceptions: the gapped-cube
  construction, and the exact re-verification of the sewn wedges, the sewn box and the sewn plate.
  \item Interval arithmetic and a compiled implementation are needed for industrial models.
  \end{itemize}
\item \textbf{Formalization.} The carrier lemma in general (contractible carriers), Waldhausen's theorem and the
existence constructions are not formalized.
\end{enumerate}




\section*{Data and code availability}
All scripts, computed results and the Lean formalization are available in the public repository~\cite{repo}.

\section*{Declaration of competing interest}
The author declares no competing interests.

\section*{Declaration of generative AI and AI-assisted technologies in the research and writing process}
During this work the author used Claude (Anthropic) to draft proofs and text, write and run the experiment scripts and
the Lean formalization, and check citations against the original sources; independent AI review passes were used to
check proofs and citations. Lean~4 with Mathlib verified the formalized lemma. The author directed the research,
reviewed and edited the content, and takes full responsibility for the publication.

\appendix
\section{A curved-triangle lemma}\label{sec:curvedtri}

Appendices~\ref{sec:proof51} and~\ref{sec:proofBprime} need maps of triangles with one curved side that reproduce a
\emph{prescribed} parametrization of that side. This is what makes maps built separately on two adjacent faces agree
on their common edge. We give an explicit construction.

\begin{lemma}\label{lem:curvedtri}
Let $\hat c\colon[t_0,t_1]\to\R^2$ be a $C^1$ regular curve, $a_1=\hat c(t_0)$, $a_2=\hat c(t_1)$, and $a_3$ a point.
\begin{itemize}
\item Put $\psi(s)=\hat c(t_0+s(t_1-t_0))-((1-s)a_1+sa_2)$ for $s\in[0,1]$.
\item In barycentric coordinates $(\lambda_1,\lambda_2,\lambda_3)$ of the reference triangle $\Delta$, put
$\mu=\lambda_1+\lambda_2$ and
\[
F(\lambda)=\lambda_1a_1+\lambda_2a_2+\lambda_3a_3+\mu^2\,\psi(\lambda_2/\mu)\qquad(F=a_3\text{ at }\mu=0).
\]
\end{itemize}
Then:
\begin{enumerate}
\item[(a)] on the side $\lambda_3=0$, $F$ is $\lambda_2\mapsto\hat c(t_0+\lambda_2(t_1-t_0))$;
\item[(b)] on the sides $\lambda_1=0$ and $\lambda_2=0$, $F$ is affine;
\item[(c)] $F$ is $C^1$ on the closed triangle $\Delta$, including at its corners.
\end{enumerate}
Moreover, consider a family of such triangles in which the chord triangles $(a_1,a_2,a_3)$ are shape-regular with
diameter $h\to0$, the speed $|\hat c'|$ is bounded above and below by positive constants (so $|t_1-t_0|=O(h)$), and
$\hat c'$ has a uniform modulus of continuity $\omega$, which we take subadditive (replacing it by its concave majorant
if necessary). Both hold in Appendices~\ref{sec:proof51} and~\ref{sec:proofBprime}, where the curves are fixed regular
$C^1$ curves on compact parameter intervals. Then for small $h$:
\begin{enumerate}
\item[(d)] $F$ is a $C^1$ diffeomorphism of $\Delta$ onto the region bounded by the arc $\hat c([t_0,t_1])$ and the
segments $[a_2,a_3]$ and $[a_3,a_1]$, and $DF=A(I+O(\omega(h)))$, where $A$ is the derivative of the affine part.
\end{enumerate}
\end{lemma}

\begin{proof}
(a) On $\lambda_3=0$ we have $\mu=1$, so $F=(1-\lambda_2)a_1+\lambda_2a_2+\psi(\lambda_2)=\hat c(t_0+\lambda_2(t_1-t_0))$.

(b) $\psi(0)=\psi(1)=0$. On $\lambda_2=0$ the correction is $\mu^2\psi(0)=0$. On $\lambda_1=0$ it is
$\mu^2\psi(1)=0$ (for $\mu>0$).

(c) For $\mu>0$, $F$ is $C^1$. Write $G=\mu^2\psi(s)$ with $s=\lambda_2/\mu$. Then
\begin{itemize}
\item $|G|\le\mu^2\max|\psi|=O(\mu^2)$, so $G$ is differentiable at $a_3$ with $DG(a_3)=0$;
\item $dG=2\mu\psi(s)\,d\mu+\mu^2\psi'(s)\,ds$ with $|ds|=O(1/\mu)$, so $|dG|=O(\mu)\to0$ as $\mu\to0$.
\end{itemize}
Hence $F$ is $C^1$ at $a_3$, with $DF(a_3)=A$. At the other two corners $\mu=1$, and $\psi$ is $C^1$ on $[0,1]$.

(d) \emph{Size of the correction.} Put $\Delta t=t_1-t_0$. Then $\psi'(s)=\Delta t\,\bigl(\hat c'(t_0+s\Delta t)-\overline{\hat
c'}\bigr)$, where $\overline{\hat c'}$ is the mean of $\hat c'$ over $[t_0,t_1]$. So $|\psi'|\le|\Delta t|\,\omega(|\Delta
t|)=O(h\,\omega(h))$, and since $\psi(0)=0$, also $|\psi|=O(h\,\omega(h))$. Hence $|DG|=O(h\,\omega(h))$ on $\Delta$, in
coordinates of the unit reference triangle.

\emph{Injectivity.} By shape regularity $\|A^{-1}\|=O(1/h)$, so $DF=A+DG=A(I+O(\omega(h)))$. For $x,y\in\Delta$ (which is
convex),
\[
|F(x)-F(y)|\ge|A(x-y)|-Ch\,\omega(h)|x-y|\ge(ch-Ch\,\omega(h))|x-y|>0
\]
for small $h$. So $F$ is injective, and its differential is invertible.

\emph{Image.} $F$ maps $\partial\Delta$ homeomorphically onto the closed curve formed by the arc and the two segments.
That curve is simple for small $h$: the arc stays within $O(h\,\omega(h))$ of its chord, and the two segments make
angles with the chord that are bounded below by shape regularity. An injective continuous map of a closed disk whose
boundary goes onto a Jordan curve has as image the closed region bounded by that curve.
\end{proof}

\emph{Remark.} The classical blending map with $\mu\,\psi(\lambda_2/\mu)$ in place of $\mu^2\psi(\lambda_2/\mu)$ is
exact on the curved side and affine on the straight sides. But it is only Lipschitz at $a_3$: its derivative there
depends on the direction of approach. The factor $\mu^2$ is what gives (c).

\emph{Standard facts used without separate checking against a source:} the Jordan curve theorem, and that an injective
continuous map of a closed disk onto a set bounded by its image boundary curve has the closed Jordan region as its
image (invariance of domain).

\section{Proof of Theorem~\ref{thm:lemma0}}\label{sec:proof51}

We use the notation and hypotheses of Section~\ref{sec:fans}; $K$ is regular.

\emph{Exit points.} For a face $f>e$, let $c(\sigma)$ be the section of $P_f$ in $\Pi(t)$, parametrized by arc length
from $b_{f,e}(t)$ and written relative to $\gamma_e(t)$. Then $(|c|^2/2)''=1+\langle c,c''\rangle\ge1-\rho_ek_f>0$ by
E2.1. So $|c|$ cannot touch the circle $\partial D_e(t)$ from inside before crossing it, and it crosses with
$\langle c',\hat r\rangle>0$, where $\hat r$ is the outward radial direction. The \emph{exit point} $k_{f,e}(t)$ is that
crossing point. It is unique by E2, and $C^1$ in $t$ by the implicit function theorem.

\emph{Margin.} $m>0$ is the minimum zone margin. $\rho_e<\tau_e$ and $W_v\subseteq\inter\Z(v)$, and by compactness all
pieces built below lie at distance at least $m$ inside the zones of their cells.

\emph{Arcs.} For a face $f$ at $v$, we use two arcs:
\begin{itemize}
\item $A_{f,v}$ is the arc of $\partial Q_f$ on $\partial W_v$, from $k_{f,e}(a_e)$ to $k_{f,e'}(a_{e'})$;
\item $\Lambda_{f,v}$ is the arc of $\lambda_v$ contributed by $f$, from $\gamma_e(a_e)$ to $\gamma_{e'}(a_{e'})$. It
consists of the two radii $[\gamma_e(a_e),k_{f,e}(a_e)]$ in the end disks, followed by $A_{f,v}$.
\end{itemize}

\subsection{The image pieces}
\begin{enumerate}
\item \emph{Vertices.} $h(v)=p_v$.
\item \emph{Edges.} On $[a_e,b_e]$, $h(e)=\gamma_e$.
\item \emph{Fans.} The fan strip is $F_{f,e}=\{\gamma_e(t)+\sigma(k_{f,e}(t)-\gamma_e(t)):t\in[a_e,b_e],\ \sigma\in[0,1]\}$.
\item \emph{Faces away from the vertices.} $Q_f=P_f\setminus\inter(\bigcup T\cup\bigcup W)$.
\item \emph{Vertex regions.} Inside $W_v$, take the cone from $p_v$ over $\lambda_v$.
  \begin{itemize}
  \item The cone over the point $\gamma_e(a_e)$ is the stub of $h(e)$.
  \item The cone over $\Lambda_{f,v}$ is the sector of $h(f)$ at $v$.
  \end{itemize}
\end{enumerate}

\subsection{Corners of $\lambda_v$}
$\lambda_v$ is piecewise $C^1$. Its corners are:
\begin{itemize}
\item the points $\gamma_e(a_e)$, where two fan radii meet;
\item the points $k_{f,e}(a_e)$, where a radius in the end disk meets the trace of $Q_f$;
\item the points where $\lambda_v$ crosses a junction curve of $\partial W_v$. By V1(b) these crossings are
transversal and avoid the corners of $\partial W_v$.
\end{itemize}
Between consecutive corners, $\lambda_v$ is a $C^1$ curve with nonvanishing tangent.

\subsection{Injectivity of the image}
\emph{Inside $T_e$.}
\begin{itemize}
\item The disks $D_e(t)$ lie in the normal planes and have radius below the reach of $A=\gamma_e([a_e,b_e])$, so they
are pairwise disjoint. Indeed, by Federer's Theorem~4.8(2),(12)~\cite[pp.~434--435]{Federer59}, if $a\in A$, $v$ lies
in the normal cone $\mathrm{Nor}(A,a)$ and $|v|<\mathrm{reach}(A)$, then $a$ is the unique nearest point of $A$ to $a+v$.
The tangent cone is the tangent line at interior points and a ray at the end points; its dual, the normal cone, is
then the normal plane, respectively a half-space, and contains the normal plane in both cases. So every point of $D_e(t)$ has nearest point $\gamma_e(t)$. This holds
whatever $\rho_e(t)$ is.
\item Inside one disk, the construction is the centre and two radii to the distinct points $k_{f,e}(t)$ and
$k_{g,e}(t)$ (E3$'$). These meet only at the centre.
\item No other piece enters $T_e$:
  \begin{itemize}
  \item non-incident pieces are excluded by S1(b)--(d) and S1(f);
  \item the patch of $f>e$ meets $D_e(t)$ only in the section arc from $b_{f,e}(t)$ to $k_{f,e}(t)$ (by the convexity
  computation above and E2.3). That arc lies in the collar of $\partial P_f$, which is removed (P0).
  \end{itemize}
\end{itemize}

\emph{Inside $W_v$.}
\begin{itemize}
\item The cone map $(s,x)\mapsto p_v+s(x-p_v)$, $s\in(0,1]$, $x\in\lambda_v$, is injective. Each ray from $p_v$ meets
$\partial W_v$ once (V1(a)), and $\lambda_v$ is simple (V1(b)).
\item No other piece enters $W_v$: S1(a), (c), (d) exclude non-incident pieces, and S1(f) excludes incident ones.
\end{itemize}

\emph{Elsewhere.}
\begin{itemize}
\item Outside the closed set $\bigcup T\cup\bigcup W$, the pieces are the embedded $Q_f$ (P0), and they are pairwise
disjoint (S1(e)).
\item On $\partial T_e$, two different $Q_f$ could meet only at exit points, which are distinct (E3$'$), or at
non-incident patches, which are excluded by S1(d).
\item On $\partial W_v$, they meet $\lambda_v$ in disjoint arcs, because $\lambda_v$ is simple.
\end{itemize}

\emph{Cell structure.}
\begin{itemize}
\item By P0, $\partial Q_f$ is one circle made of the curves $k_{f,e}$ and the arcs $A_{f,v}$.
\item By V1(b) there is exactly one arc per corner of $f$. So the pieces form a 2-regular graph isomorphic to
$\partial f$; its only cycle is $\partial Q_f$, and the cyclic order is forced.
\item The pieces meet exactly as follows:
  \begin{itemize}
  \item $Q_f\cap F_{f,e}=k_{f,e}$, by E2.3 and S1(f);
  \item $(Q_f\cup\text{fans})\cap\text{sector}_{f,v}=\Lambda_{f,v}$, by strict star-shapedness and S1(f).
  \end{itemize}
\item Gluing two disks along a boundary arc that is their whole intersection gives a disk. By induction, the union
$h(f)$ is a closed disk with $\partial h(f)=\bigcup_{e<f}h(e)$.
\end{itemize}

\emph{Zones.}
\begin{itemize}
\item $Q_f\subseteq P_f\subseteq\Z(f)$.
\item $F_{f,e}\subseteq T_e\subseteq N_\rho(\gamma_e)\subseteq\Z(e)\subseteq\Z(f)$.
\item The cone parts lie in $W_v\subseteq\Z(v)$, which is contained in $\Z(e)$ and $\Z(f)$ for incident $e$ and $f$.
\item $h(\inter e)\subseteq\gamma_e\cup W_v\subseteq\Z(e)$.
\item All inclusions hold with margin at least $m$.
\end{itemize}

\subsection{The map $h$}
We build a polyhedral domain and a homeomorphism onto the union of the pieces that is $C^1$, with injective
differential, on each closed domain cell (except at the apexes), and that agrees on shared edges.

\emph{Domain.} Let $L$ be a polyhedral subdivision of $|K|$ with 2-cells for the fans, for the pieces $Q_f$ and for the
sectors, in the incidence pattern of the image pieces. Each sector is subdivided into \emph{sub-cones}: cones from the
vertex $v$ over one segment of a polygonal arc of the link of $v$. That arc has a vertex at the preimage of every corner
of $\lambda_v$, so the image of a sub-cone is the cone over a $C^1$ piece of $\lambda_v$.

\emph{The parametrization $\phi$.} Fix once and for all a parametrization $\phi$ of $\lambda_v$ by the polygonal link,
$C^1$ with $\phi'\ne0$ on each link segment, such that:
\begin{itemize}
\item on each end radius $[\gamma_e(a_e),k_{f,e}(a_e)]$ it is affine in the fan parameter $\sigma$;
\item on each arc $A_{f,v}$ it is some fixed $C^1$ regular parametrization.
\end{itemize}
The maps below all use this $\phi$.

\emph{Fans.} Use the chart $(t,\sigma)\mapsto\gamma_e(t)+\sigma(k_{f,e}(t)-\gamma_e(t))$ on $[a_e,b_e]\times[0,1]$,
precomposed with a bilinear map from the quadrilateral domain cell onto this rectangle.
\begin{itemize}
\item It is $C^1$ (since $k_{f,e}$ is).
\item Its differential is injective: $\partial_\sigma=k-\gamma\in\Pi(t)$ is nonzero, and $\partial_t$ has
$\gamma'$-component at least $1-\sigma\kappa_\gamma\rho_e>0$.
\item It fixes the parametrization of the crease $k_{f,e}$ by $t$, and that of the end radius at $t=a_e$ by $\sigma$.
\end{itemize}

\emph{Sectors.} On each sub-cone, map $v+s(y-v)$ to $p_v+s(\phi(y)-p_v)$. This map is $C^1$ up to the boundary
away from $v$, with injective differential, because $\phi'\ne0$ and $\phi'$ is transverse to the rays (V1(a)). It agrees
with the fan on the end radii (both affine in $\sigma$), and with the map on $Q_f$ along $A_{f,v}$ (both use $\phi$).

\emph{$Q_f$.} In the parameter domain of $P_f$, $\partial Q_f$ is a piecewise $C^1$ curve. Its pieces are $k_{f,e}$,
parametrized by $t$, and $A_{f,v}$, parametrized by $\phi$. Its corners are at the points $k_{f,e}(a_e)$ and at the
junction crossings.
\begin{itemize}
\item Subdivide $Q_f$ into curved triangles, each with at most one curved edge. The curved edges are pieces of
$\partial Q_f$ with their prescribed parametrizations, and the corners of $\partial Q_f$ are vertices. Triangulate the
domain cell of $Q_f$ in $L$ to match.
\item On each curved triangle, use the map of Lemma~\ref{lem:curvedtri}. It reproduces the prescribed
parametrization on the curved edge, is affine on the straight edges, is $C^1$ up to the corners, and is a diffeomorphism
for a fine enough subdivision. As in Appendix~\ref{sec:proofBprime}, put one interior vertex near each corner of
$\partial Q_f$ so that no triangle has two curved edges.
\item Composed with the chart of $P_f$, this gives maps that are $C^1$ up to the boundary, have injective
differential, and agree on shared edges.
\end{itemize}

These cell maps agree on shared edges, so together they define a homeomorphism $h\colon|L|=|K|\to\bigcup$ pieces. It
sends each cell of $K$ onto the corresponding image cell.

\subsection{The PL step}
\emph{Metric.} We measure Lipschitz constants on each shell of $|L|$ in its intrinsic length metric. That metric is
bi-Lipschitz equivalent to the extrinsic one, and in it a Lipschitz bound on each cell is a global bound. Injectivity
between shells is argued separately at the end.

\emph{$h$ is bi-Lipschitz.} $h$ is injective on a compact space, so it suffices that it is locally bi-Lipschitz.
\begin{itemize}
\item \emph{Inside a cell}, this follows from the injective differential.
\item \emph{Along the creases}, the two half-planes of the tangent cone must differ.
  \begin{itemize}
  \item Along $k$ (fan and $Q_f$): in $\Pi(t)$ the fan leaves $k$ in direction $-\hat r$ and the patch section in
  direction $c'$ with $\langle c',\hat r\rangle>0$. Since $k'\notin\Pi(t)$, the two half-planes could coincide only if
  $c'$ were parallel to $-\hat r$, which is excluded.
  \item Along $\gamma_e$ (two fans): E3$'$ and compactness.
  \item Along $\partial W_v$ (cone and outside piece): the cone's half-tangent $p_v-x$ points strictly into $W_v$ (strict
  star-shapedness, V1(a)), and the outside piece leaves strictly outward (transversality, V1(b)). For a fan, the
  tangent plane along the end radius contains $\partial_t=(1-\sigma)\gamma'+\sigma k'$, which does not lie in $\Pi(a_e)$.
  \item Along the rays from $p_v$ to a corner of $\lambda_v$ (between two sub-cones): by homogeneity, the opening angle
  along the ray is constant away from $p_v$, and it is positive because the radial projection of $\lambda_v$ to the unit
  sphere about $p_v$ has no fold at the corner (see the corner arguments below and V1(a)).
  \end{itemize}
\item \emph{At $k(a_e)$} (fan, $Q_f$ and two cone sub-sectors). Here $\rho_e'(a_e)=0$, so $k'\perp\hat r$. Let
$Q=\{\langle\cdot,\gamma'\rangle\le0,\langle\cdot,\hat r\rangle\le0\}$ and $w=p_v-k(a_e)$; then $w\in\inter Q$ by strict
star-shapedness.
  \begin{itemize}
  \item The fan sector $\mathrm{span}^+(k',-\hat r)$ lies in $\{\langle\cdot,\gamma'\rangle\ge0\}$.
  \item The $Q_f$ half-plane, bounded by $\R k'$ and containing $c'$, lies in $\{\langle\cdot,\hat r\rangle\ge0\}$.
  \item The two cone sub-sectors $\mathrm{span}^+(-\hat r,w)$ and $\mathrm{span}^+(w,-k')$ lie in $Q$.
  \end{itemize}
  So the fan and $Q_f$ sectors lie in the closure of the complement of $Q$, and the cone sectors lie in $Q$. They share
  only the rays $k'$ (fan and $Q_f$), $-\hat r$ (fan and cone), $w$ (between the two cone sub-sectors) and $-k'$ (cone
  and $Q_f$). The two cone sub-sectors do not fold: if $-\hat r$, $w$, $-k'$ were coplanar, then $w\in\inter Q$ would
  force $w=-a\hat r-bk'$ with $a,b>0$, strictly between them. Hence the tangent cone is an embedded disk.
\item \emph{At $\gamma(a_e)$} (two fans and the adjacent cone sub-cones). The fans lie in
$\{\langle\cdot,\gamma'\rangle\ge0\}$, and the cone sub-cones in $\{\langle\cdot,\gamma'\rangle\le0\}$ with
$\langle p_v-\gamma(a_e),\gamma'\rangle<0$. They share only the rays $\hat r_f\ne\hat r_g$ (E3$'$). The two cone
sub-cones, $\mathrm{span}^+(\hat r_f,w)$ and $\mathrm{span}^+(w,\hat r_g)$ with $w=p_v-\gamma(a_e)$, do not fold onto each
other, because $\hat r_f\ne\hat r_g$ and both are orthogonal to $\gamma'$ while $w$ is not.
\item \emph{At a junction crossing $x$ of $\lambda_v$.} The tangent cone $T_x$ of $W_v$ at $x$ is the union of two
half-spaces, since $W_v$ is a union of two convex sets near $x$. Let $\ell_1,\ell_2$ be the half-tangents of $\lambda_v$
at $x$, one in each face of $T_x$.
  \begin{itemize}
  \item The cone sub-cones near $x$ consist of segments from $p_v$ to points of $\lambda_v$. By strict star-shapedness
  these segments run into $W_v$, so the sub-cones lie in $W_v$ near $x$. Their tangent sectors
  $\mathrm{span}^+(\ell_1,w)$ and $\mathrm{span}^+(w,\ell_2)$, with $w=p_v-x$, therefore lie in $T_x$. (This uses
  star-shapedness itself, not the convexity of $T_x$, which fails.)
  \item The two sub-cones do not fold: projected along $w$, the half-tangents $\ell_1,\ell_2$ point in different
  directions, because the radial projection of $\lambda_v$ is locally injective (each ray meets $\partial W_v$ once,
  V1(a)) and $\lambda_v$ is transverse to the rays.
  \item By the transversality in V1(b), the tangent sector of the outside piece lies in the closure of the complement
  of $T_x$.
  \end{itemize}
  So the sectors meet only along $\ell_1$, $w$ and $\ell_2$, and the tangent cone is an embedded disk.
\item \emph{At the apex $p_v$.} $h$ is homogeneous of degree one on each sector. The radial projection of $\lambda_v$
to the unit sphere about $p_v$ is injective (V1(a)). It is locally bi-Lipschitz: by transversality between corners,
and by the no-fold arguments above at the corners. $\phi$ is bi-Lipschitz on each link segment. Hence $h$ is
bi-Lipschitz on the star of $v$.
\end{itemize}

\emph{Secant approximation.} Take a triangulation of $|L|$ with these properties:
\begin{itemize}
\item it refines every cell of $L$ and has every corner as a vertex;
\item it is shape-regular off the sector cells;
\item it is radial on the sector cells: its triangles there have the apex $v$ as a vertex and a sub-segment of one link
segment as the opposite side.
\end{itemize}
The two kinds of triangles meet only along those base segments. Let $S$ be the secant map, linear on each triangle and
equal to $h$ at the vertices.
\begin{itemize}
\item \emph{Off the sectors.} $h$ is $C^1$ up to the boundary on each closed triangle. The standard per-triangle
estimate gives $\operatorname{Lip}(S-h)\le C\,\omega(Dh,\mathrm{mesh})/\sin\theta_{\min}$, where $\omega$ is a modulus
of continuity and $\theta_{\min}$ is bounded below by shape regularity.
\item \emph{On a radial triangle $(v,y_0,y_1)$.} Its base lies on one link segment. Write $y=(1-u)y_0+uy_1$ with
$u\in[0,1]$, let $\ell=|y_1-y_0|$, and let $\tau$ be the unit vector of the base; derivatives of $e$ below are with
respect to arc length $\ell u$.
  \begin{itemize}
  \item A point $x=v+s(y-v)$ has barycentric coordinates $(1-s,\,s(1-u),\,su)$. So $S(x)=p_v+s(\hat\phi(y)-p_v)$,
  where $\hat\phi$ is the linear interpolation of $\phi$ on the base, and
  \[
  (S-h)(x)=s\,e(y),\qquad e=\hat\phi-\phi .
  \]
  This is exactly homogeneous of degree one in $s$.
  \item Differentiating: $\partial_s x=y-v$ gives $\partial_s(S-h)=e(y)$. And $\partial_u x=s\ell\tau$ while
  $\partial_u(S-h)=s\ell e'$, so $D(S-h)\cdot\tau=e'$, with no division by the distance to the apex.
  \item Hence $\|D(S-h)\|\le(|e|/|y-v|+|e'|)/\sin\beta$, where $\beta$ is the angle between the ray $y-v$ and the base.
  $\beta$ is bounded below by the fixed domain geometry, since the link is transverse to the rays.
  \item Since $\hat\phi'$ is the mean of $\phi'$ over the base, $|e'|\le\omega(\phi',\mathrm{mesh})$ and
  $|e|\le\mathrm{mesh}\cdot\omega(\phi',\mathrm{mesh})$.
  \item $D(S-h)$ is constant along rays, so the bound holds uniformly up to the apex, and $S-h$ is continuous at $v$.
  So $\operatorname{Lip}(S-h)\le C\,\omega(\phi',\mathrm{mesh})\to0$.
  \end{itemize}
  This uses only that $\phi$ is $C^1$ on each link segment; no curvature of $\lambda_v$ is needed.
\item \emph{Conclusion.} With the intrinsic metric, $\operatorname{Lip}(S-h)\to0$ and $|S-h|\to0$ globally. For a fine
enough triangulation:
  \begin{itemize}
  \item $\operatorname{Lip}(S-h)$ is below the bi-Lipschitz constant of $h$, so $S$ is injective on each shell, by
  $|S(x)-S(y)|\ge|h(x)-h(y)|-\operatorname{Lip}(S-h)\,|x-y|$;
  \item $h$ maps different shells to disjoint compact sets at some distance $\delta>0$, so once $|S-h|<\delta/2$,
  $S$ maps them to disjoint sets;
  \item $|S-h|<m$, so $S(\inter\sigma)\subseteq\Z(\sigma)$.
  \end{itemize}
  $S$ is PL, so it is a realization. $\square$
\end{itemize}

\emph{Standard facts used without separate checking against a source:} that the tangent cone of a $C^1$ regular arc
is its tangent line (a ray at the end points); the implicit function theorem; gluing disks along boundary arcs; the per-triangle
secant-approximation estimate for $C^1$ maps; bi-Lipschitz equivalence of the intrinsic and extrinsic metrics on a
compact polyhedron; that a map whose tangent cones are embedded disks with positive opening angles is locally
bi-Lipschitz; and that an injective, locally bi-Lipschitz map of a compact space is bi-Lipschitz.

\section{Proof of Proposition~\ref{prop:Bprime}}\label{sec:proofBprime}

The idea is to take the re-trimmed surface $\Sigma$ itself as the model of $|K|$, and the identity as a topological
realization, then replace it by a fine secant approximation. Notation is that of Proposition~\ref{prop:Bprime}.
\begin{itemize}
\item $X_f\colon U_f\to\R^3$ are the face parametrizations, and $D_f\subset U_f$ the compact parameter regions.
\item $\phi\colon|K|\to\Sigma$ is the homeomorphism of C.
\end{itemize}

\subsection{Margin}
By C, $\phi$ maps each closed cell of $K$ onto the corresponding closed cell of $\Sigma$. By B1--B3 each of these lies in
the zone of its cell. Closed cells are compact and finitely many, and zones are open. So there is $m>0$ such that the
$m$-neighbourhood of every closed cell of $\Sigma$ lies in its zone.

\subsection{A curved triangulation of $\Sigma$}
Fix a small $h>0$.
\begin{enumerate}
\item \emph{Edge vertices.} On each curve $\tilde\gamma_e$, fix one $C^2$ regular parametrization (by arc length, say),
and choose vertices on it at 3D spacing about $h$, including its end points. A closed edge gets at least three
vertices. These choices are made once, in $\R^3$, and are shared by all faces containing $e$ (both sides of a seam).
\item \emph{Cutting along the edge graph.} In each $D_f$, pull back the edge vertices by $X_f$; this is possible because
$X_f$ is injective on $D_f$. Cut $D_f$ along its edge graph $G_f$ (R), seams included. Note that $\partial D_f\subseteq G_f$: a point of
$\partial D_f$ off $G_f$ would be mapped by $X_f$ into the open face, which by C is an open surface homeomorphic, through
$X_f^{-1}$, to an open subset of $U_f$; by invariance of domain such a point is interior to $D_f$. The result of cutting
is a finite union of compact regions, each bounded by arcs of $G_f$ (a seam arc bounds a region on each of its two
sides). A region need not be simply connected: faces with inner wires give multiply connected regions. Join consecutive
edge vertices along each arc by chords. For small $h$ each boundary component of a region becomes a simple closed
polygon, these polygons are disjoint, and they bound a polygonal region close to the original one, by R: the arcs are
simple $C^2$ regular, and the sector angles lie in $(0,2\pi)$.
\item \emph{Triangulation.} Triangulate each chord region shape-regularly, with no extra vertices on its boundary. At
every vertex of $G_f$, including smooth passages such as seam vertices, put one interior vertex near it in each adjacent
sector and join it to the vertex, so that no triangle has two curved edges. The shape-regularity constant depends on
the smallest sector angle.
\item \emph{Curved triangles.} Replace each chord by its arc. An arc deviates from its chord by $o(h)\cdot h$, which is
much smaller than $h\sin\theta_{\min}$, so the curved triangles tile $D_f$, and every arc of $G_f$, seams included, is a
union of edges of the triangulation.
\end{enumerate}
Let $T$ be the resulting abstract complex. For small $h$ it is a simplicial complex, and it subdivides the cells of
$\Sigma$, hence, through $\phi$, the cells of $K$. We metrize $|T|$ by standard (equilateral) simplices and use the intrinsic
length metric on each shell.

\emph{The map $h_0\colon|T|\to\Sigma$.}
\begin{itemize}
\item On a triangle with one curved edge, use Lemma~\ref{lem:curvedtri}, with the curved edge parametrized by the
pull-back of the fixed parametrization of $\tilde\gamma_e$. Precompose with the affine map from the standard simplex,
and compose with $X_f$.
\item On a straight triangle, use the affine map, composed with $X_f$.
\end{itemize}
Then:
\begin{itemize}
\item \emph{$h_0$ agrees on shared edges.}
  \begin{itemize}
  \item An interior straight edge is affine in $U_f$ on both sides.
  \item An edge on $\tilde\gamma_e$ between faces $f$ and $g$ is $s\mapsto\tilde\gamma_e(t_0+s(t_1-t_0))$ from both
  sides, by Lemma~\ref{lem:curvedtri}(a).
  \item A seam edge lies inside $D_f$, and both sides use the same pull-back of the fixed parametrization of the
  seam curve, so it is mapped the same way from both sides.
  \end{itemize}
\item \emph{$h_0$ is a homeomorphism.} It is a homeomorphism onto $\Sigma$, by Lemma~\ref{lem:curvedtri}(d) and the
injectivity of $X_f$ on $D_f$.
\item \emph{Regularity.} On each closed triangle $h_0$ is $C^1$ up to the boundary, with injective differential. This
follows from Lemma~\ref{lem:curvedtri}(c),(d), since $X_f$ is $C^2$ regular on a neighbourhood of $D_f$.
\end{itemize}

\subsection{$h_0$ is bi-Lipschitz on each shell}
On a shell, $|T|$ is compact and $h_0$ is injective, so it suffices that $h_0$ is locally bi-Lipschitz.
\begin{itemize}
\item \emph{Inside a face,} including across interior edges and at interior vertices of $T$: the map to $D_f$ is a
piecewise-$C^1$ homeomorphism with invertible differential on each closed piece. Its tangent map at a vertex is a
sectorwise linear, orientation-preserving map of the plane of degree one, hence a homeomorphism. Composing with the
regular parametrization $X_f$ keeps it locally bi-Lipschitz.
\item \emph{Along an edge (seams included), away from its end points:} by T the two inward tangent half-planes are
distinct. For a seam they are opposite in the tangent plane of $S_f$. The angle between them is continuous and
positive, hence bounded away from $0$ on compact subarcs; near the end points the vertex argument below takes over. So the tangent cone at
each point is a union of two half-planes that meet only along their common line. This also covers vertices of $T$ that
lie on $\tilde\gamma_e$: there the triangle sectors on each side tile the two half-planes.
\item \emph{At a vertex $c_v$:} by T the tangent cone of $\Sigma$, the union of the tangent sectors of the edge graphs at
the preimages of $c_v$, is an embedded disk; by R each sector has an opening angle in $(0,2\pi)$. $h_0$ is $C^1$ up to $c_v$ on every
closed triangle, so the blow-up argument applies: rescale at $c_v$, pass to the tangent cone, and use $C^1$ convergence
on each sector.
\end{itemize}

\subsection{Secant approximation}
Refine $T$ by affine subdivision of each simplex; this keeps shape regularity and keeps $h_0$ $C^1$ on each original
closed triangle. Let $S$ be the secant map: linear on each refined triangle and equal to $h_0$ at the vertices.
\begin{itemize}
\item \emph{The estimate.} On each triangle, $\operatorname{Lip}(S-h_0)\le C\,\omega(Dh_0,\mathrm{mesh})/\sin\theta_{\min}$ and
$|S-h_0|\le\mathrm{mesh}\cdot\omega(Dh_0,\mathrm{mesh})$, where $\omega$ is a modulus of continuity. In the intrinsic
metric of each shell these bounds are global.
\item \emph{Within a shell.} Once $\operatorname{Lip}(S-h_0)$ is below the bi-Lipschitz constant of $h_0$, $S$ is
injective there.
\item \emph{Between shells.} $h_0$ maps different shells to disjoint compact sets at distance $\delta>0$. Once
$|S-h_0|<\delta/2$, $S$ maps them to disjoint sets.
\item \emph{Zones.} Once $|S-h_0|<m$: every point of a closed cell lies on a simplex of $T$ inside that cell, so
$S(\bar\sigma)\subseteq\Z(\sigma)$ for vertices, edges and faces. In particular $S(\inter\sigma)\subseteq\Z(\sigma)$.
\end{itemize}
$S\circ h_0^{-1}\circ\phi$ is then a PL embedding of $|K|$, linear on the simplices of a triangulation that subdivides the
cells of $K$, with every cell in its zone. That is a realization. $\square$

\emph{Standard facts used without separate checking against a source:}
\begin{itemize}
\item existence of shape-regular triangulations of a polygonal region, possibly multiply connected, with prescribed
boundary vertices and corner fans as described;
\item the per-triangle secant-approximation estimate;
\item that a sectorwise linear, orientation-preserving map of the plane of degree one is a homeomorphism;
\item that a map that is $C^1$ up to the vertex on each sector, with an embedded tangent cone of positive opening
angles, is locally bi-Lipschitz (blow-up);
\item that an injective, locally bi-Lipschitz map of a compact space is bi-Lipschitz.
\end{itemize}

\section{Formal verification in Lean}\label{sec:lean}

The following are proved in Lean~4 with Mathlib (repository directory \texttt{lean/}). The build has no
\texttt{sorry}, and the results depend only on the axioms \texttt{propext}, \texttt{Classical.choice} and
\texttt{Quot.sound}.

\begin{center}\small
\begin{tabular}{p{6.6cm}p{6.4cm}}
\toprule
Lean statement & Content \\
\midrule
\texttt{inner\_eq\_zero\_of\_conical} & at a nearest point $x$, $\langle p-x,x-a\rangle=0$ (Lemma~\ref{lem:cone}) \\
\texttt{inner\_sub\_eq\_norm\_sq\_of\_conical} & $\langle p-x,p-a\rangle=|p-x|^2$ \\
\texttt{not\_isCritical\_of\_conical} & $p$ is not critical (Lemma~\ref{lem:cone}) \\
\texttt{not\_isCritical\_of\_conicalCover} & pointwise form of Corollary~\ref{cor:lambda} \\
\texttt{carrier\_mono}, \texttt{carrier\_eq\_zone\_of\_mono} & carriers are monotone, and equal the zones when the
zones are augmented \\
\texttt{carried\_of\_realization}, \texttt{image\_closedCell\_subset\_carrier} & realizations are carried \\
\texttt{carried\_straightLine}, \texttt{continuous\_straightLine} & Lemma~\ref{lem:carrier} when the carriers
are convex \\
\bottomrule
\end{tabular}
\end{center}

The statements are for an arbitrary real inner product space. Criticality is formalized as
$0\in\conv\{p-x:x\text{ nearest}\}$.

\emph{How the Lean statements correspond to the paper.}
\begin{itemize}
\item The hypothesis \texttt{LocallyConicalAt M a r} is radial closure of $M$ inside the open ball $B(a,r)$. The
paper's definition of $\rho(a)$ implies it for $r<\rho(a)$; the two are not the same statement.
\item Lean's \texttt{Carried} is pointwise on open cells. By \texttt{carrier\_mono} it agrees with the paper's
closed-cell definition.
\item Positivity of $\lambda(M)$ for polyhedra (Corollary~\ref{cor:lambda}) is proved on paper only.
\end{itemize}

\begin{thebibliography}{99}

\bibitem{APS00} L.-E. Andersson, T.~J. Peters, N.~F. Stewart, Equivalence of topological form for curvilinear
geometric objects, \emph{Int.\ J.\ Comput.\ Geom.\ Appl.} 10(6) (2000) 609--622.

\bibitem{ASZ07} L.-E. Andersson, N.~F. Stewart, M. Zidani, Error analysis for operations in solid modeling in the
presence of uncertainty, \emph{SIAM J.\ Sci.\ Comput.} 29(2) (2007) 811--826.

\bibitem{ASZ06} L.-E. Andersson, N.~F. Stewart, M. Zidani, Conditions for use of a non-selfintersection conjecture,
\emph{Comput.\ Aided Geom.\ Design} 23(7) (2006) 599--611.

\bibitem{CCS05} F. Chazal, D. Cohen-Steiner, A condition for isotopic approximation, \emph{Graphical Models} 67(5)
(2005) 390--404.

\bibitem{CCSL09dcg} F. Chazal, D. Cohen-Steiner, A. Lieutier, A sampling theory for compact sets in Euclidean space,
\emph{Discrete Comput.\ Geom.} 41(3) (2009) 461--479.

\bibitem{CCSL09cg} F. Chazal, D. Cohen-Steiner, A. Lieutier, Normal cone approximation and offset shape isotopy,
\emph{Comput.\ Geom.} 42(6--7) (2009) 566--581.

\bibitem{ChazalLieutier07} F. Chazal, A. Lieutier, Stability and computation of topological invariants of solids
in $\mathbb R^n$, \emph{Discrete Comput.\ Geom.} 37 (2007) 601--617.

\bibitem{dREG22} S. di Rocco, D. Eklund, O. Gäfvert, Computing geometric feature sizes for algebraic manifolds,
arXiv:2209.01654.

\bibitem{Durfee83} A.~H. Durfee, Neighborhoods of algebraic sets, \emph{Trans.\ Amer.\ Math.\ Soc.} 276(2) (1983) 517--530.

\bibitem{Federer59} H. Federer, Curvature measures, \emph{Trans.\ Amer.\ Math.\ Soc.} 93(3) (1959) 418--491.

\bibitem{Fu85} J.~H.~G. Fu, Tubular neighborhoods in Euclidean spaces, \emph{Duke Math.\ J.} 52(4) (1985) 1025--1046.

\bibitem{Jackson95} D.~J. Jackson, Boundary representation modelling with local tolerances, in \emph{Proc.\ 3rd ACM
Symp.\ Solid Modeling and Applications (SMA '95)}, 1995, 247--254.

\bibitem{Milenkovic88} V.~J. Milenkovic, Verifiable implementations of geometric algorithms using finite precision
arithmetic, \emph{Artificial Intelligence} 37(1--3) (1988) 377--401.


\bibitem{Moise77} E.~E. Moise, \emph{Geometric Topology in Dimensions 2 and 3}, Graduate Texts in Mathematics 47,
Springer, New York, 1977.

\bibitem{OCCT} Open Cascade, \emph{OpenCascade Technology documentation}, version 8.0.1: BRep format specification,
\S\S7.2--7.4 (shape tolerances).

\bibitem{QiShapiro06} J. Qi, V. Shapiro, $\varepsilon$-topological formulation of tolerant solid modeling,
\emph{Computer-Aided Design} 38(4) (2006) 367--377.

\bibitem{QiShapiro06b} J. Qi, V. Shapiro, Geometric interoperability with epsilon solidity, \emph{J.\ Comput.\ Inf.\
Sci.\ Eng.} 6(3) (2006) 213--220.

\bibitem{SSP01} T. Sakkalis, G. Shen, N.~M. Patrikalakis, Topological and geometric properties of interval solid
models, \emph{Graphical Models} 63(3) (2001) 163--175.

\bibitem{SPB04} T. Sakkalis, T.~J. Peters, J. Bisceglio, Isotopic approximations and interval solids,
\emph{Computer-Aided Design} 36 (2004) 1089--1100.

\bibitem{Segal90} M. Segal, Using tolerances to guarantee valid polyhedral modeling results, \emph{Computer
Graphics (Proc.\ SIGGRAPH '90)} 24(4) (1990) 105--114.

\bibitem{Shapiro08} V. Shapiro, Homotopy conditions for tolerant geometric queries, in \emph{Reliable Implementation
of Real Number Algorithms}, LNCS 5045, Springer, 2008, 162--180.

\bibitem{SZ08} N.~F. Stewart, M. Zidani, Transfinite interpolation for well-definition in error analysis in solid
modelling, in \emph{Reliable Implementation of Real Number Algorithms}, LNCS 5045, Springer, 2008, 181--192.

\bibitem{Waldhausen68} F. Waldhausen, On irreducible 3-manifolds which are sufficiently large, \emph{Ann.\ of
Math.} 87 (1968) 56--88.

\bibitem{repo} G. Zhang, \texttt{tolerant-solids} (code, proofs, data), \url{https://github.com/eddiezgq/tolerant-solids}
(public repository).

\end{thebibliography}

\end{document}
